\documentclass[11pt,a4paper]{article}
\usepackage[T1]{fontenc}
\usepackage[utf8]{inputenc}
\usepackage{lmodern,microtype}
\usepackage[margin=27mm,headheight=14pt]{geometry}
\usepackage{amsmath,amssymb,amsthm,mathtools}
\usepackage{bm}
\usepackage{booktabs,longtable,array,enumitem}
\usepackage{tikz}
\usepackage{xcolor}
\usepackage{hyperref}
\usepackage{fancyhdr}
\hypersetup{colorlinks=true,linkcolor=blue!45!black,citecolor=blue!45!black,urlcolor=blue!45!black,
 pdftitle={A Pole Hierarchy for Valley-Monotone Compositions},
 pdfauthor={Research draft prepared for Vladimir Reshetnikov}}
\pagestyle{fancy}\fancyhf{}\fancyhead[L]{\small Valley-monotone compositions}\fancyhead[R]{\small Research draft}\fancyfoot[C]{\thepage}
\setlength{\parskip}{4pt}\setlength{\parindent}{14pt}
\setlength{\emergencystretch}{2em}
\numberwithin{equation}{section}
\newtheorem{theorem}{Theorem}[section]
\newtheorem{lemma}[theorem]{Lemma}
\newtheorem{proposition}[theorem]{Proposition}
\newtheorem{corollary}[theorem]{Corollary}
\theoremstyle{definition}\newtheorem{definition}[theorem]{Definition}
\theoremstyle{remark}\newtheorem{remark}[theorem]{Remark}
\newcommand{\N}{\mathbb N}
\newcommand{\C}{\mathbb C}
\newcommand{\E}{\mathbb E}
\newcommand{\Prob}{\mathbb P}
\newcommand{\Var}{\operatorname{Var}}
\newcommand{\Res}{\operatorname{Res}}
\newcommand{\coeff}[2]{[#1^{#2}]}
\newcommand{\dd}{\mathrm d}
\newcommand{\ee}{\mathrm e}
\newcommand{\calP}{\mathcal P}
\newcommand{\calD}{\mathcal D}
\newcommand{\calU}{\mathcal U}
\DeclareMathOperator{\Log}{Log}
\newcommand{\Cov}{\operatorname{Cov}}
\DeclareMathOperator{\Dir}{Dirichlet}
\DeclareMathOperator{\Unif}{Uniform}
\newcommand{\wrt}[1]{\textup{[}\textsf{write:}~#1\textup{]}}
\newenvironment{writenote}[1][]{\par\smallskip\begingroup\small
 \noindent\textbf{[write]}\if\relax\detokenize{#1}\relax\else~\textbf{#1.}\fi\ }%
 {\par\endgroup\smallskip}
\title{\textbf{A Pole Hierarchy for\\Valley-Monotone Compositions}\\[7pt]
\large Resolving a bargraph asymptotic conjecture linked to OEIS A001523,\\
with exponential expansions, limit laws, and Lambert--\(W\) inversion}
\author{Research draft prepared for Vladimir Reshetnikov}
\date{October 1, 2026}
\begin{document}
\maketitle
\begin{abstract}
A composition is valley-monotone when the heights of its interior valley
plateaux are weakly increasing from left to right. These are the
non-decreasing bargraphs of Fl\'orez, Ram\'irez, and Villamizar (2024).
Their area-counting sequence was conjectured to have purely exponential
growth, on the basis of a ratio at area~100. We give a self-contained
analytic proof of that growth law, exact definitions of its constants,
and integer-arithmetic certificates for their numerical values.
The growth constant is determined directly by the generating function
of the classical weakly unimodal composition sequence OEIS A001523.

The proof reveals a stronger structure. A decomposition by minimum part
produces infinitely many genuine, simple positive poles, increasing to
one, whose coefficient contributions alternate in sign. Consequently the
area generating function is not D-finite and its coefficients are not
P-recursive. We prove a finite-pole expansion to every prescribed radius
below one, a central limit theorem for the number of bottom-level valleys,
and an exponentially tight limiting terminal component. The accumulating
real poles admit an all-orders inverse-logarithmic expansion expressed
through Lambert~\(W\), with an error beyond every inverse-logarithmic order.
We also give exponential-sector inversion and the necessary integer
threshold qualifications. The claims are conventional mathematical
proofs, accompanied by reproducible checks, not Lean formalizations or
independently peer-reviewed results.
\end{abstract}
\begin{writenote}[Added 5 October 2026, batch~98]
Since batch~98 the report has two Parts. Part~I, summarized above, is the
original report. Part~II (Section~\ref{vmb:cw:sec:front} onward) adds
manuscript~04 of batch~98: it answers Part~I's extreme-height question with
a lattice-Gumbel law for the maximum part, introduces a valley-weighted
model with a first-order transition at an explicit critical weight, and
proves critical-window, Dirichlet and Dirichlet-mixture laws, an exact
height-tail identity with an expansion of the height-cutoff pole, and
inverses of these expansions. Its labels carry the prefix \texttt{vmb:cw:}.
\end{writenote}
\tableofcontents
\newpage
\phantomsection\addcontentsline{toc}{part}{Part I.\texorpdfstring{\quad}{ } A pole hierarchy for valley-monotone compositions}

\section{The question, the OEIS connection, and the scope of the result}

\subsection{A published asymptotic conjecture}
Let \(d(n)\) count nonempty compositions of \(n\) whose valley heights are
weakly increasing. A valley is a maximal constant run with larger parts
on both sides; endpoint runs are not valleys. For example,
\((3,1,2,2,4,2,3)\) has valley heights \(1,2\), whereas
\((3,2,3,1,2)\) has heights \(2,1\) and is excluded.
The area of the corresponding bargraph is the sum of its column heights.
The sequence begins
\[
1,2,4,8,16,32,64,128,256,512,1023,2042,4071,8106,16121,\ldots.
\]
The convention throughout is \(d(0)=0\).

Conjecture~2.3 of Fl\'orez, Ram\'irez, and Villamizar~\cite{FRV}
proposes an exponential equivalent for this sequence. The displayed
numerical estimates in that conjecture are
\begin{equation}\label{vmb:eq:published}
\lambda_{100}=1.9764206780767492582\ldots,
\qquad C_{100}=0.6039015148966519\ldots.
\end{equation}
Here the growth estimate is the finite ratio \(d(100)/d(99)\), not an
exact definition of the limiting exponential base. We establish
\begin{equation}\label{vmb:eq:main-intro}
d(n)=C\lambda^n\bigl(1+O(\theta^n)\bigr),\qquad 0<\theta<1,
\end{equation}
with
\begin{align*}
\lambda&=1.97642067794583573531465245217332\ldots,\\
C&=0.60390151906189910496030721536293\ldots.
\end{align*}
Thus the proposed \emph{form} is correct, while some displayed numerical
digits must be replaced. Interpreting a finite decimal base literally in
an asymptotic equivalent would be incorrect: an arbitrarily small fixed
error in that base becomes exponentially significant as \(n\) grows.

\subsection{What is, and is not, being claimed about OEIS}
Let
\begin{equation}\label{vmb:eq:A}
A(q)=\sum_{n\geq0}a(n)q^n,\qquad a(n)=\text{OEIS A001523}(n).
\end{equation}
This is the classical sequence of weakly unimodal compositions, with
\(a(0)=1\). Its leading stretched-exponential asymptotic is already
known; the OEIS entry credits Auluck (1951) and links the 2024 bargraph
paper~\cite{OEIS,FRV}. We do \emph{not} claim to solve a previously open
leading-asymptotic problem for A001523 itself.

Instead, the classical sequence is the exact input to a new pole
analysis of the associated valley-monotone class. In particular, the
reciprocal growth constant \(\rho=\lambda^{-1}\) is the unique solution
in \((0,1)\) of
\begin{equation}\label{vmb:eq:OEISroot}
\boxed{\quad \rho(1-\rho)A(\rho)=1+\rho.\quad}
\end{equation}
No OEIS identifier for \(d(n)\) was verified in the searches for this
report, so none is invented. The connection to OEIS is both explicit and
structural, rather than an accidental match of a few terms.

\subsection{Relationship to ProveIt and novelty boundary}
The repository's \texttt{Analysis/Transseries} project separates
asymptotic scales, flatness, and inverse expansions from the Fabius
function development~\cite{ProveIt}. The present article fits that
program, but does not require any unproved theorem from the repository.
A repository search for \texttt{bargraphs} and \texttt{A001523} returned
no matching files. This was a search-level overlap check, not an audit of
every archive in the repository. The inspected repository tree had SHA
\texttt{29aca108ed25d714f31b1316512777fbbdc8e006}.

The generating-function decomposition is due to the 2024 paper; it is
reproved below in minimum-part notation. The contributions developed
here are the analytic resolution of the conjectured exponential form,
the full positive-pole hierarchy and non-D-finiteness, the pole
accumulation expansion, and the probabilistic and inverse consequences.
Targeted searches located no earlier proof of this particular conjecture
or these stronger conclusions. That is a bounded literature-search
statement, not a guarantee of priority. General meromorphic coefficient
extraction and analytic inversion are standard tools~\cite{FS}.

\begin{writenote}[Added 2 October 2026, batch~75]
This report is a single manuscript, printed in full: manuscript~04 of
batch~75 of ProveIt's incoming-reports intake, delivered as the archive
\texttt{OEIS\_\allowbreak Valley\_\allowbreak Monotone\_\allowbreak Bargraphs.zip}
in commit \texttt{4b874cea0} and placed in commit \texttt{6ea60e367} in the
research-report collection, among the enumerative-combinatorics reports
(\nolinkurl{SetTheory/Cardinals/docs/reports}), not in the
\texttt{Analysis/Transseries} project named above. The identifier called a
``tree SHA'' above is the ProveIt \emph{commit} \texttt{29aca108e}
(1~October 2026), an ancestor of the placement commit; it is this
manuscript's pin. Nothing in this report is formalized in Lean or Rocq, and
placement in the collection confers no formal status on any statement here.
No OEIS number for the sequence $d(n)$ was verified, and the writing step
assigns none; A001523 is the input sequence, not the target. The writing step
changed the labels (every label now carries the prefix \texttt{vmb:}),
renamed Theorem~\ref{vmb:thm:W} (delivered as ``Pole accumulation
transseries''), set the bibliography ragged-right, and added these
\textsf{[write]} notes and one reference \cite{vmb-tai}; no other statement,
proof or number of the manuscript was altered. The Lambert-$W$ and
rounding steps of Sections~\ref{vmb:sec:accumulation}
and~\ref{vmb:sec:inverse} are instances of the transseries-and-inversion
volume \cite{vmb-tai}, as the notes there say; no transseries in that
volume's sense is constructed here.
\end{writenote}

\begin{writenote}[Added 5 October 2026, batch~98]
The note above describes Part~I, which was a single manuscript printed in
full. Since batch~98 the report also has a Part~II, built from a second
manuscript (manuscript~04 of batch~98, pin \texttt{70ab31a1e}, placed in
commit \texttt{9353a7171}); Section~\ref{vmb:cw:sec:provenance} records its
provenance and Section~\ref{vmb:cw:sec:notation} its notation, which differs
from Part~I's in places listed there.
\end{writenote}

\section{A decomposition by minimum part}\label{vmb:sec:decomposition}

\subsection{Classes and notation}
For an integer \(j\geq1\), let \(\calU_j\) be the nonempty weakly
unimodal compositions whose parts are at least \(j\), and let
\(\calD_j\) be the nonempty valley-monotone compositions with that same
lower bound. Write
\begin{equation}
U_j(q)=\sum_{\alpha\in\calU_j}q^{|\alpha|},\qquad
D_j(q)=\sum_{\alpha\in\calD_j}q^{|\alpha|},\qquad
d_j(n)=\coeff qn D_j(q).
\end{equation}
Here \(|\alpha|\) is the sum of the parts. In particular,
\(U_1=A-1\), \(D_1(q)=\sum_{n\geq1}d(n)q^n\).
Define
\begin{equation}\label{vmb:eq:HKQ}
H_j(q)=\frac{q^j}{1-q^j},\quad
K_j(q)=H_j(q)U_{j+1}(q),\quad
Q_j(q)=1-q^j-q^jU_{j+1}(q).
\end{equation}
Thus
\begin{equation}\label{vmb:eq:QK}
Q_j(q)=(1-q^j)(1-K_j(q)).
\end{equation}
The series \(H_j\) counts a nonempty run of parts equal to~\(j\).

\begin{lemma}[Unimodal decomposition]\label{vmb:lem:U}
As formal power series,
\begin{equation}\label{vmb:eq:Urec}
U_j=H_j+\frac{U_{j+1}}{(1-q^j)^2}.
\end{equation}
Equivalently,
\begin{equation}\label{vmb:eq:Usum}
U_j(q)=\sum_{h\geq j}\frac{q^h}{1-q^h}
                   \prod_{r=j}^{h-1}(1-q^r)^{-2}.
\end{equation}
\end{lemma}
\begin{proof}
An object either consists only of parts~\(j\), or has a nonempty middle
object with all parts greater than~\(j\). In the second case, unimodality
allows parts~\(j\) only in the initial and final runs. Those runs may be
empty and are chosen independently, giving \((1-q^j)^{-2}\).
This proves the recurrence. Iteration classifies an object by its maximum
part~\(h\) and gives the sum. For each coefficient of fixed degree, only
finitely many values of \(h\) contribute, so the formal iteration is
legitimate.
\end{proof}

\begin{lemma}[Valley-monotone decomposition]\label{vmb:lem:D}
As formal power series,
\begin{equation}\label{vmb:eq:Drec}
\boxed{\quad D_j=H_j+
 \frac{D_{j+1}}{(1-q^j)^2(1-K_j)}
 =H_j+\frac{D_{j+1}}{(1-q^j)Q_j}.\quad}
\end{equation}
\end{lemma}
\begin{proof}
The all-\(j\) object contributes \(H_j\). Otherwise split the
composition into its maximal blocks of parts greater than~\(j\), with
runs of \(j\)'s between them and possibly at the two ends.
Each internal separating run is a valley of height~\(j\).

Every block except the last must be unimodal. Indeed, an internal valley
in an earlier block would have height greater than~\(j\) and would be
followed by a separating valley of height~\(j\), violating the required
order. A finite sequence without an internal valley is weakly unimodal:
a downward change followed later by an upward change necessarily creates
an intervening valley. Conversely, unimodal earlier blocks create no
additional valleys and the last block may be any member of
\(\calD_{j+1}\).

The unique decomposition is therefore
\[
(\text{initial }j\text{-run})\;
(\calU_{j+1}\,\text{nonempty }j\text{-run})^*\;
\calD_{j+1}\;
(\text{final }j\text{-run}).
\]
The two endpoint factors contribute \((1-q^j)^{-2}\); the repeated unit
has generating function \(K_j\). This proves the formula, including
uniqueness and the absence of overcounting.
\end{proof}

\begin{corollary}[Iterated series]\label{vmb:cor:Dsum}
For \(j\geq1\),
\begin{equation}\label{vmb:eq:Dsum}
D_j(q)=\sum_{h\geq j} H_h(q)
\prod_{r=j}^{h-1}\frac{1}{(1-q^r)Q_r(q)}.
\end{equation}
In particular,
\begin{equation}\label{vmb:eq:paperGF}
D_1(q)=\sum_{h\geq1}
\frac{q^h}{\displaystyle\prod_{r=1}^{h}(1-q^r)
                 \prod_{r=1}^{h-1}Q_r(q)}.
\end{equation}
\end{corollary}
This agrees with the area specialization in~\cite{FRV}, where the
paper's \(B_j\) is our \(U_{j+1}\). The boundary condition at large~\(j\)
is coefficientwise disappearance: \(U_j,D_j\) have no terms below degree
\(j\). It uniquely determines both recurrences.

\subsection{The renewal kernel in terms of A001523}
From \eqref{vmb:eq:Urec} at \(j=1\),
\[
U_2=(1-q)^2U_1-q(1-q).
\]
Consequently,
\begin{equation}\label{vmb:eq:K-A}
K_1(q)=q(1-q)A(q)-q
      =\sum_{n\geq3}\bigl(a(n-1)-a(n-2)\bigr)q^n.
\end{equation}
Thus the differences of a classical OEIS sequence count the repeated
units. This identity will also provide the probability distribution
underlying the limit theorem.

\section{Analytic continuation and the positive-pole hierarchy}

\subsection{Normal convergence inside the unit disk}
\begin{lemma}[Uniform estimates]\label{vmb:lem:bounds}
Fix \(0<R<1\). Uniformly for \(|q|\leq R\), all the functions \(U_j\)
are analytic and
\begin{equation}\label{vmb:eq:Ubound}
|U_j(q)|\leq B_R R^j,
\qquad |K_j(q)|=O_R(R^{2j+1}),
\qquad Q_j(q)=1+O_R(R^j),
\end{equation}
where \(B_R\) is finite and independent of~\(j\).
\end{lemma}
\begin{proof}
The infinite product
\(P_R=\prod_{r\geq1}(1-R^r)^{-2}\) is finite, since the sum of the
absolute logarithms converges. Each term in \eqref{vmb:eq:Usum} is bounded
by \(P_R R^h/(1-R)\). Summing for \(h\geq j\) gives
\(B_R=P_R/(1-R)^2\), and local uniform convergence proves analyticity.
The other estimates follow from \eqref{vmb:eq:HKQ}.
\end{proof}

\begin{theorem}[Meromorphic continuation]\label{vmb:thm:meromorphic}
Every \(D_j\) extends to a single-valued meromorphic function on
\(|q|<1\). Its possible poles there are among the zeros of
\(Q_r\), \(r\geq j\). On each compact subdisk there are only finitely
many possible poles.
\end{theorem}
\begin{proof}
For sufficiently large \(r\), uniformly on a chosen compact subdisk,
\(Q_r\) stays away from zero and
\[
\frac{1}{(1-q^r)Q_r(q)}=1+O_R(R^r).
\]
The corresponding tail products in \eqref{vmb:eq:Dsum} are uniformly
bounded, because \(\sum_rR^r\) converges. The summands are therefore
\(O_R(R^h)\) after separating finitely many initial denominator factors.
The tail sum converges normally off those factors' zeros. Multiplying by
the finite product of initial \(Q_r\)'s makes the sum analytic locally.
Each initial \(Q_r\) is analytic and not identically zero, so its zeros
are isolated and finite in a compact subdisk. These observations prove
all the assertions.
\end{proof}

\subsection{Distinguished real roots}
\begin{lemma}\label{vmb:lem:roots}
For every \(j\geq1\) there is a unique \(\rho_j\in(0,1)\) with
\(K_j(\rho_j)=1\), equivalently \(Q_j(\rho_j)=0\). Moreover,
\begin{equation}\label{vmb:eq:rholadder}
0<\rho_1<\rho_2<\cdots<1,
\qquad \rho_j\longrightarrow1,
\qquad Q_j'(\rho_j)<0.
\end{equation}
The only zero of \(Q_j\) on \(|q|\leq\rho_j\) is \(q=\rho_j\).
\end{lemma}
\begin{proof}
On \((0,1)\), \(K_j\) is strictly increasing, starts at zero, and tends
to infinity. The last claim follows already from the subclass consisting
of one part greater than~\(j\), together with the nonempty separating
run. Thus \(\rho_j\) exists uniquely and its derivative is positive.
Equation~\eqref{vmb:eq:QK} then gives
\(Q_j'(\rho_j)=-(1-\rho_j^j)K_j'(\rho_j)<0\).

For positive \(q<1\), both \(H_j>H_{j+1}\) and
\(U_{j+1}>U_{j+2}\). Hence \(K_j>K_{j+1}\), proving strict ordering
of their unit crossings. The estimates in Lemma~\ref{vmb:lem:bounds} show
that \(K_j(q)\to0\) for every fixed \(q<1\); this forces
\(\rho_j\to1\).

If \(|q|<\rho_j\), positivity gives
\(|K_j(q)|\leq K_j(|q|)<1\), excluding zeros. On the boundary,
equality with \(K_j(q)=1\) in the triangle inequality would force all
nonzero summands of \(K_j(q)\) to be nonnegative real. Its support
contains both \(2j+1\) and \(2j+2\): use a separating part~\(j\) and
a singleton of height \(j+1\) or \(j+2\). Their consecutive exponents
force \(q/\rho_j=1\). This proves boundary uniqueness.
\end{proof}

\begin{theorem}[Genuine simple poles and alternating amplitudes]
\label{vmb:thm:ladder}
For every \(k\geq j\), \(D_j\) has a genuine simple pole at
\(\rho_k\). In particular, write the principal part of \(D_1\) there
as
\[
\frac{C_k}{1-q/\rho_k}.
\]
Then
\begin{equation}\label{vmb:eq:Ck}
C_k=
\frac{D_{k+1}(\rho_k)}{(1-\rho_k^k)(-\rho_k Q_k'(\rho_k))}
\prod_{r=1}^{k-1}\frac{1}{(1-\rho_k^r)Q_r(\rho_k)},
\qquad \operatorname{sgn}(C_k)=(-1)^{k-1}.
\end{equation}
These are exactly the positive real poles inside the unit disk.
\end{theorem}
\begin{proof}
By Lemma~\ref{vmb:lem:roots}, \(D_{k+1}\) is analytic at \(\rho_k\):
all of its possible pole factors have indices greater than~\(k\) and
larger minimum zero modulus. Its defining positive series converges
there and has positive value. Applying \eqref{vmb:eq:Drec} at level~\(k\)
therefore produces a genuine simple pole, of amplitude
\[
\frac{D_{k+1}(\rho_k)}{(1-\rho_k^k)(-\rho_k Q_k'(\rho_k))}>0.
\]
For every \(r<k\), the recurrence at level~\(r\) adds an analytic
term and multiplies the inherited pole by a finite nonzero factor
\(((1-q^r)Q_r(q))^{-1}\). Since \(Q_r(\rho_k)<0\), each such factor
changes its sign. This proves the formula, noncancellation, and sign.
Theorem~\ref{vmb:thm:meromorphic} and uniqueness of each positive zero
exclude any other positive poles.
\end{proof}

\begin{corollary}[Non-D-finiteness and non-P-recursiveness]
\label{vmb:cor:nonD}
None of the functions \(D_j\) is D-finite over \(\C(q)\). In
particular \(d(n)\) satisfies no eventual nonzero linear recurrence
with polynomial coefficients in~\(n\).
\end{corollary}
\begin{proof}
A function satisfying a linear differential equation with polynomial
coefficients can have finite-plane singularities only at the finitely
many zeros of its leading coefficient. At every other point, standard
local existence for a linear analytic differential equation continues
its analytic branches holomorphically. The infinitely many poles
\(\rho_k\), \(k\geq j\), contradict this requirement.

For completeness, a polynomial-coefficient recurrence for the Taylor
coefficients, after multiplication by powers of \(q\) and summation,
yields a linear differential equation for the generating function, using
\(q\,\dd/\dd q\) to encode multiplication by~\(n\). Any finite
initial discrepancy is polynomial and may be annihilated by further
differentiation. Thus an eventual recurrence would also imply
D-finiteness, which has just been excluded.
\end{proof}

\section{The exponential law and certified constants}\label{vmb:sec:leading}

\begin{theorem}[Leading asymptotics at every fixed minimum part]
\label{vmb:thm:leading}
For each fixed \(j\geq1\), there is \(R_j>\rho_j\), with
\(R_j<1\), such that
\begin{equation}\label{vmb:eq:leadingj}
d_j(n)=c_j\rho_j^{-n}+O_j(R_j^{-n}),
\qquad
c_j=\frac{D_{j+1}(\rho_j)}
{(1-\rho_j^j)^2\rho_jK_j'(\rho_j)}>0.
\end{equation}
For \(j=1\), \(c_1=C_1=C\) and \(\rho_1=\rho\), so this proves
\eqref{vmb:eq:main-intro} and the intended form of Conjecture~2.3 of~\cite{FRV}.
\end{theorem}
\begin{proof}
Lemma~\ref{vmb:lem:roots} and Theorem~\ref{vmb:thm:ladder} give a unique simple
pole on the circle of convergence. Local finiteness of all possible
poles permits a slightly larger closed disk containing no others.
Subtract the principal part \(c_j/(1-q/\rho_j)\). The difference is
analytic on a neighborhood of that disk. Cauchy's coefficient bound on
its boundary gives the stated error. The formula for \(c_j\) follows
from \eqref{vmb:eq:QK}. Equation~\eqref{vmb:eq:K-A} gives
\eqref{vmb:eq:OEISroot}.
\end{proof}

There are no nonzero inverse-power corrections multiplying the leading
exponential: for every fixed \(M\),
\begin{equation}
\frac{d(n)}{C\lambda^n}=1+o(n^{-M}).
\end{equation}
The missing information is instead in genuinely smaller exponential
sectors. This is why an inverse-power fit to the leading term would
conceal the analytic structure.

\subsection{Exact definitions and integer-arithmetic enclosures}
The constants are defined by
\begin{equation}\label{vmb:eq:constants}
K_1(\rho)=1,
\qquad \lambda=\rho^{-1},
\qquad C=\frac{D_2(\rho)}{(1-\rho)^2\rho K_1'(\rho)}.
\end{equation}
The accompanying \texttt{verify.py} proves the following enclosures with
outward-rounded rational intervals, using integers only:
\begin{equation}\label{vmb:eq:rhobound}
\begin{gathered}
0.50596515770080666537292224451833314155847753935777<\rho,\\
\rho<0.50596515770080666537292224451833314155847753935778.
\end{gathered}
\end{equation}
\begin{gather}
1.976420677945835735314652452172
<\lambda<1.976420677945835735314652452175,
\label{vmb:eq:lambdabound}\\
0.603901519061899104960307215362
<C<0.603901519061899104960307215364.
\label{vmb:eq:Cbound}
\end{gather}
The narrower machine-produced intervals are retained in
\texttt{verification.json}. These bounds certify, in particular, the
numerical distinction between the limit and the finite-area estimates.

\begin{proposition}[Tail bounds used in the certificate]\label{vmb:prop:tail}
For \(0<q<1\), put
\(S_J(q)=q^{J+1}/(1-q)\). If \(S_J(q)<1\), then
\begin{equation}\label{vmb:eq:tails}
0\leq U_{J+1}(q),D_{J+1}(q)\leq\frac{S_J(q)}{1-S_J(q)},
\quad
0\leq U_{J+1}'(q)\leq\frac{S_J'(q)}{(1-S_J(q))^2}.
\end{equation}
\end{proposition}
\begin{proof}
\(S_J\) counts a single allowed part. All nonempty ordered sequences
of such parts have generating function \(S_J/(1-S_J)\), and contain
both restricted classes as subclasses. Coefficientwise domination proves
the first two inequalities and, after differentiating nonnegative
coefficients, the derivative bound.
\end{proof}
The program uses \(J=250\) and endpoints in \(10^{-90}\mathbb Z\),
propagates these tail intervals down to level~2, and tests the sign of
\(Q_1\) at the two rational endpoints in \eqref{vmb:eq:rhobound}.
To enclose \(C\), it uses the equivalent expression
\begin{equation}\label{vmb:eq:Ccertificate}
C=\frac{D_2(\rho)}
{\rho(1-\rho)\{1+U_2(\rho)+\rho U_2'(\rho)\}}.
\end{equation}
Every division checks that its denominator interval excludes zero.
No unbounded numerical tail is discarded.

\section{Exponential expansions to arbitrary subunit radius}

\begin{theorem}[Finite-pole expansion]\label{vmb:thm:sectors}
Choose \(0<R<1\) such that \(D_1\) has no pole on \(|q|=R\), and let
\(\calP_R\) be its finite set of actual poles inside that circle. Write
the principal part at \(\zeta\in\calP_R\) in the form
\[
\sum_{\ell=1}^{m_\zeta}
\frac{c_{\zeta,\ell}}{(1-q/\zeta)^\ell}.
\]
Then
\begin{equation}\label{vmb:eq:allsectors}
\boxed{\quad d(n)=
\sum_{\zeta\in\calP_R}\sum_{\ell=1}^{m_\zeta}
 c_{\zeta,\ell}\binom{n+\ell-1}{\ell-1}\zeta^{-n}
 +O_R(R^{-n}).\quad}
\end{equation}
\end{theorem}
\begin{proof}
Subtract all the displayed principal parts. The result is analytic in a
neighborhood of the closed disk, so its coefficients are \(O_R(R^{-n})\)
by Cauchy's formula. The coefficient of each subtracted rational term is
the displayed binomial coefficient times \(\zeta^{-n}\).
\end{proof}

For positive real poles, \(m_{\rho_k}=1\) and
\(c_{\rho_k,1}=C_k\). Thus these sectors have no polynomial
prefactors. Nonreal conjugate pairs may yield exponentially small
oscillations; multiplicities or coincidences among complex zeros may
yield polynomial factors. Formula~\eqref{vmb:eq:allsectors} makes no
unsupported simplicity assumption for those zeros.

The numerical positive-real hierarchy begins as follows. The table's
ordering is by \emph{positive real pole}, not by modulus of all complex
poles. Apart from the first row, these decimal entries are numerical
checks rather than rational certificates.
\begin{center}
\begin{tabular}{@{}rrrr@{}}
\toprule
\(k\)&\(\rho_k\)&\(\rho_k^{-1}\)&\(C_k\)\\
\midrule
1&0.505965157700807&1.976420677945836&\phantom{-}0.603901519061899\\
2&0.628520593778803&1.591037763755333&-0.434266071916591\\
3&0.694915251028261&1.439024396889127&\phantom{-}0.097127120171435\\
4&0.738032724347822&1.354953468877238&-0.008189474598318\\
5&0.768802380990748&1.300724379536014&\phantom{-}0.000290088602322\\
6&0.792097830048713&1.262470318771736&-0.000004651415177\\
\bottomrule
\end{tabular}
\end{center}

\begin{remark}[Meaning of ``all orders'']
For every fixed exponential cutoff \(R<1\), only finitely many sectors
are required and the remainder is proved. This does not assert that the
infinite sum of all pole contributions converges, that the unit circle
is a natural boundary, or that the positive-real poles alone constitute
the complete expansion. In particular, we have not certified that
\(\rho_2\) is the second-smallest modulus among \emph{all} poles.
An all-orders algebraic expansion around the dominant term is trivial
here; the finite-pole theorem is the substantive exponential refinement.
\end{remark}

\section{A renewal limit law and a tight terminal component}

\subsection{Bottom-level valleys}
Choose uniformly from the \(d(n)\) objects of area~\(n\), and let
\(V_n\) be the number of valleys of height~1. Introduce a variable~\(u\)
marking these valleys. The decomposition in Lemma~\ref{vmb:lem:D} gives
\begin{equation}\label{vmb:eq:marked}
D(q,u)=H_1(q)+\frac{B(q)}{1-uK(q)},
\qquad B(q)=\frac{D_2(q)}{(1-q)^2},\quad K(q)=K_1(q).
\end{equation}
The repeated unit contributes exactly one height-1 valley.
Define a positive integer random variable \(Z\) by
\begin{equation}\label{vmb:eq:Z}
\Prob(Z=s)=\coeff qs K(q)\,\rho^s.
\end{equation}
This is a probability distribution because \(K(\rho)=1\). It has
exponential moments near zero and includes both \(s=3\) and \(s=4\)
in its support. Put
\begin{equation}\label{vmb:eq:MV}
M=\E Z=\rho K'(\rho),\qquad
V=\Var Z=\rho^2K''(\rho)+\rho K'(\rho)-M^2>0.
\end{equation}

\begin{theorem}[Gaussian bottom-valley count]\label{vmb:thm:CLT}
As \(n\to\infty\),
\begin{equation}\label{vmb:eq:CLT}
\frac{V_n-\mu n}{\sqrt{\sigma^2n}}
\ \Longrightarrow\ \mathcal N(0,1),
\qquad \mu=\frac1M,\qquad \sigma^2=\frac{V}{M^3}.
\end{equation}
Moreover \(\E V_n=\mu n+O(1)\) and
\(\Var V_n=\sigma^2n+O(1)\). Numerical evaluation gives
\[
\mu=0.1372487560663833542\ldots,
\qquad \sigma^2=0.0368978521690088765\ldots.
\]
\end{theorem}
\begin{proof}
The implicit function theorem gives an analytic root \(r(u)\) near
\(u=1\), satisfying \(uK(r(u))=1\) and \(r(1)=\rho\).
Choose a closed disk slightly larger than \(\rho\) in which
\eqref{vmb:eq:marked} at \(u=1\) has only its simple pole. Compactness and
analytic perturbation preserve exactly this one nearby pole for complex
\(u\) sufficiently close to~1. Subtracting it uniformly gives
\begin{equation}\label{vmb:eq:quasipower}
\E[u^{V_n}]=\frac{C(u)}{C(1)}
\left(\frac{\rho}{r(u)}\right)^n(1+O(\vartheta^n)),
\quad
C(u)=\frac{B(r(u))}{u r(u)K'(r(u))},
\end{equation}
for some \(\vartheta<1\), locally uniformly in~\(u\).

Set \(t(s)=\log r(e^s)\). Taking logarithms of the root equation gives
\(\log K(e^{t(s)})=-s\). The first two derivatives of
\(\log K(e^t)\) at \(t=\log\rho\) are respectively \(M\) and \(V\).
Therefore
\[
t'(0)=-M^{-1},\qquad t''(0)=-VM^{-3}.
\]
Taking logarithms in \eqref{vmb:eq:quasipower} and differentiating twice
proves the mean and variance statements. Substituting
\(s=iz/\sqrt{\sigma^2n}\), subtracting the centering term, and using
the analytic Taylor expansion gives a limiting characteristic function
\(\exp(-z^2/2)\). This proves the asserted weak convergence.
\end{proof}

\subsection{Why almost all valleys are at height one}
In the nonflat decomposition, let \(T_n\in\calD_2\) be the final
block with parts greater than one, and let \(L_n,R_n\geq0\) be the
lengths of the two endpoint runs of ones. The all-ones composition has
probability \(1/d(n)\). On this exceptional object, define
\((L_n,R_n)=(0,0)\) and let \(T_n=\dagger\) be a cemetery object of
area zero. This convention is asymptotically irrelevant.

\begin{theorem}[Boltzmann terminal limit]\label{vmb:thm:terminal}
The triple \((L_n,T_n,R_n)\) converges in distribution to independent
components \((L,T,R)\), where
\begin{equation}\label{vmb:eq:terminal}
\Prob(L=\ell)=\Prob(R=\ell)=(1-\rho)\rho^\ell,
\qquad
\Prob(T=\tau)=\frac{\rho^{|\tau|}}{D_2(\rho)}
\quad(\tau\in\calD_2).
\end{equation}
The sizes of these components have uniformly bounded exponential moments
for all sufficiently large~\(n\). The number of valleys above height
one is therefore \(O_{\Prob}(1)\), and the total number of valleys has
the same central limit theorem as \(V_n\).
\end{theorem}
\begin{proof}
For a fixed triple \((\ell,\tau,r)\), its contribution is
\(q^{\ell+|\tau|+r}/(1-K(q))\). Comparing its dominant residue with
that of \eqref{vmb:eq:marked} at \(u=1\) gives limiting probability
\[
\frac{\rho^{\ell+|\tau|+r}}{B(\rho)}
=(1-\rho)^2\rho^{\ell+r}
 \frac{\rho^{|\tau|}}{D_2(\rho)},
\]
which factors as stated and sums to one.

To justify tightness and exponential moments, mark endpoint length by
\(v\) and terminal area by~\(w\). The numerator becomes
\(D_2(wq)/(1-vq)^2\). Since \(D_2\) is analytic past \(\rho\),
this numerator remains analytic on a common larger disk when real
\(v,w>1\) are sufficiently close to~1. The same residue argument
then bounds the marked expectation uniformly in~\(n\).
All valleys above height one lie in~\(T_n\), and their number is at
most its area. Their contribution divided by \(\sqrt n\) tends to
zero in probability, so Slutsky's theorem transfers the central limit
law to the total valley count.
\end{proof}

The limiting terminal area has mean
\[
\E|T|=\frac{\rho D_2'(\rho)}{D_2(\rho)}
=6.2446300331882976716\ldots.
\]
This decimal is numerical, not one of the certified intervals above.
The same arguments work at every fixed level~\(j\), with \(K_j\),
\(D_{j+1}\), and runs of parts~\(j\). As an additional consequence of
Theorem~\ref{vmb:thm:leading},
\begin{equation}
\Prob(\text{all parts are at least }2)
\sim\frac{c_2}{c_1}\left(\frac{\rho_1}{\rho_2}\right)^n.
\end{equation}
Here \(c_2>0\) is the \emph{dominant amplitude of \(D_2\)}, not the
negative amplitude \(C_2\) of the second positive pole of~\(D_1\).

\section{How the real poles accumulate at one}\label{vmb:sec:accumulation}

\subsection{An exact product comparison}
For \(0<q<1\), define
\begin{equation}
P_{j+1}(q)=\prod_{h\geq j+1}(1-q^h)^{-2}.
\end{equation}
\begin{lemma}\label{vmb:lem:productcomparison}
If \(y=q^j\), then
\begin{equation}\label{vmb:eq:Pcomparison}
\frac{1-yq}{2-yq}(P_{j+1}-1)
\ \leq U_{j+1}\ \leq\frac12(P_{j+1}-1).
\end{equation}
\end{lemma}
\begin{proof}
Let \(A_h=\prod_{r=j+1}^h(1-q^r)^{-2}\), with \(A_j=1\).
The term of \(U_{j+1}\) with maximum part~\(h\) is
\[
\frac{q^h}{1-q^h}A_{h-1}
=\frac{1-q^h}{2-q^h}(A_h-A_{h-1}).
\]
The multiplicative factor lies between
\((1-yq)/(2-yq)\) and~\(1/2\). Sum the positive telescoping
increments and pass to the limit.
\end{proof}

Put
\begin{equation}\label{vmb:eq:xt}
t_j=-\log\rho_j,\qquad x_j=jt_j,\qquad y_j=e^{-x_j}=\rho_j^j,
\qquad c=\log2.
\end{equation}
At the pole, the defining equation is
\(U_{j+1}(\rho_j)=(1-y_j)/y_j\).

\begin{lemma}[Pole-scale equation]\label{vmb:lem:scale}
As \(j\to\infty\), \(x_j\to\infty\), \(x_j\sim\log j\), and
\begin{equation}\label{vmb:eq:scale}
\frac{2j e^{-x_j}}{x_j}
=x_j+c+O\!\left(\frac{x_j^3}{j}\right).
\end{equation}
\end{lemma}
\begin{proof}
We know \(t_j\to0\). We first show \(y_j\to0\). Otherwise some
subsequence has \(y_j\geq\delta>0\). The singleton subclass gives
\(U_{j+1}(\rho_j)\geq\rho_j^{j+1}/(1-\rho_j)\to\infty\),
whereas \((1-y_j)/y_j\) is bounded, a contradiction.
Thus \(x_j\to\infty\).

At the root, \eqref{vmb:eq:Pcomparison} gives
\begin{equation}\label{vmb:eq:logP}
\log P_{j+1}(\rho_j)=x_j+c+O(y_j).
\end{equation}
On the other hand, expanding the product logarithm absolutely gives,
with \(t=t_j\), \(y=y_j\),
\begin{align}
\log P_{j+1}(e^{-t})
 &=2\sum_{k\geq1}\frac{y^k}{k(e^{kt}-1)}\notag\\
 &=\frac{2y}{t}+O\!\left(y+\frac{y^2}{t}\right).
\label{vmb:eq:logPexp}
\end{align}
For the last estimate use
\((e^t-1)^{-1}=t^{-1}+O(1)\) for \(k=1\), and
\((e^{kt}-1)^{-1}\leq(kt)^{-1}\) for \(k\geq2\).

Since \(y\to0\), the error is \(o(y/t)\), while the right side of
\eqref{vmb:eq:logP} tends to infinity. It follows that
\(2y/t\sim x\), or \(2j e^{-x}\sim x^2\). Taking logarithms gives
\(x+2\log x=\log(2j)+o(1)\), hence \(x\sim\log j\) and
\(y\asymp x^2/j\). Substituting this estimate back into
\eqref{vmb:eq:logPexp} and \eqref{vmb:eq:logP} gives \eqref{vmb:eq:scale}.
\end{proof}

\subsection{A Lambert--\texorpdfstring{\(W\)}{W} expansion to all logarithmic orders}
\begin{theorem}[Pole accumulation: a Lambert--\(W\) expansion]\label{vmb:thm:W}
Let \(X_j>0\) be the unique positive solution of
\begin{equation}\label{vmb:eq:X}
e^{X_j}X_j(X_j+c)=2j,\qquad c=\log2,
\end{equation}
and set
\begin{equation}\label{vmb:eq:w}
w_j=2W\!\left(\sqrt{j/2}\right),
\end{equation}
where \(W\) is the positive real Lambert function. Then
\begin{equation}\label{vmb:eq:xX}
x_j=X_j+O\!\left(\frac{(\log j)^2}{j}\right).
\end{equation}
For every fixed \(m\geq1\),
\begin{equation}\label{vmb:eq:xW}
x_j=w_j+\sum_{k=1}^{m}\frac{\alpha_k(c)}{w_j^k}
 +O_m(w_j^{-m-1})+O\!\left(\frac{w_j^2}{j}\right),
\end{equation}
where
\begin{align}
\alpha_1(c)&=-c,\notag\\
\alpha_2(c)&=\frac{c(c+4)}2,\notag\\
\alpha_3(c)&=-\frac{c(c^2+6c+12)}3,\notag\\
\alpha_4(c)&=\frac{c(3c+8)(c^2+6c+12)}{12}.
\label{vmb:eq:alphas}
\end{align}
The full series for \(X_j-w_j\) is convergent for sufficiently large
\(w_j\); the difference between \(x_j\) and this logarithmic sector
is smaller than every fixed inverse power of~\(\log j\).
\end{theorem}
\begin{proof}
Let \(F_j(x)=2je^{-x}/x-x-c\). Its derivative is negative.
Lemma~\ref{vmb:lem:scale} gives
\(F_j(x_j)=O(x_j^3/j)\), while \(F_j(X_j)=0\).
First, \(F_j'\leq-1\) gives
\(|x_j-X_j|=O(x_j^3/j)=o(1)\). Thus throughout the interval between
these two points, \(|F_j'|\asymp x_j\), because
\(2je^{-x_j}/x_j\sim x_j\). A second application of the mean value
theorem improves the difference to \(O(x_j^2/j)\), proving
\eqref{vmb:eq:xX}.

The definition of \(w=w_j\) is equivalent to \(e^w w^2=2j\).
Write \(X=w+\varepsilon\), \(z=1/w\). Taking logarithms of
\eqref{vmb:eq:X} and cancelling the equation for~\(w\) gives
\begin{equation}\label{vmb:eq:implicit-eps}
\varepsilon+2\log(1+z\varepsilon)
 +\log\left(1+\frac{cz}{1+z\varepsilon}\right)=0.
\end{equation}
This is analytic near \((z,\varepsilon)=(0,0)\), with derivative~1
with respect to~\(\varepsilon\). The analytic implicit function
theorem supplies a convergent series
\(\varepsilon=\sum_{k\geq1}\alpha_k(c)z^k\).
Coefficient comparison gives \eqref{vmb:eq:alphas}, and gives every further
coefficient recursively: substitute the already computed terms, and
negate the coefficient of \(z^k\) to obtain \(\alpha_k\). The new
coefficient occurs only in the first, linear term at that order.
Combining this convergent expansion with \eqref{vmb:eq:xX} proves all claims.
\end{proof}

In particular,
\begin{equation}\label{vmb:eq:rhorough}
1-\rho_j\sim\frac{\log j}{j},\qquad
x_j=\log j-2\log\log j+\log2+o(1).
\end{equation}
Higher expansions for \(\rho_j=e^{-x_j/j}\) and the growth bases
\(\rho_j^{-1}=e^{x_j/j}\) now follow by ordinary exponential
composition, with the error propagated explicitly. The theorem is a
large-\(j\) statement about the pole locations; it does not by itself
make the fixed-\(j\) coefficient estimate uniform when \(j\) grows with
\(n\).

\begin{writenote}[Added 2 October 2026, batch~75: an instance of repository results]
The manuscript titled Theorem~\ref{vmb:thm:W} ``Pole accumulation
transseries''; it is a single-scale Lambert-$W$ expansion of the pole
positions, and the theorem has been renamed accordingly. Taking logarithms,
\eqref{vmb:eq:X} reads $X+2\log X=\log(2j)-\log(1+c/X)$. Its dominant
block $X+2\log X=\log(2j)$ is solved exactly by the Lambert core
\texttt{p0:thm:lambert-core} of the transseries-and-inversion volume
\cite{vmb-tai} (there $a=1$, $b=2>0$, principal branch), which gives
$w_j=2W(\sqrt{j/2})$. The correction equation \eqref{vmb:eq:implicit-eps}
is that volume's reversion about an arbitrary core,
\texttt{p0:thm:core-reversion}, with $\Lambda=1$, $h=0$ and small
parameter $t=z=1/w_j$; its triangular recurrence is the coefficient
recursion described in the proof, and \texttt{code/expansion.py} implements
it. The leading equation of Corollary~\ref{vmb:cor:polecount}, in the
variable $x+c$, is another instance of that core (there $a=b=1$), which
gives $x+c=W(4/t)$. These are instances of that apparatus, and no novelty
is claimed for the method. What is specific to this report is the
pole-scale equation (Lemma~\ref{vmb:lem:scale}) and the error terms in
\eqref{vmb:eq:xX} and~\eqref{vmb:eq:xW}.
\end{writenote}

\subsection{Inverting the real-pole counting function}
\begin{corollary}\label{vmb:cor:polecount}
Let \(N_+(r)=\#\{j\geq1:\rho_j\leq r\}\), and put \(t=-\log r\).
Then, as \(r\uparrow1\),
\begin{equation}\label{vmb:eq:polecount}
N_+(r)=\frac{W(4/t)-\log2}{t}
       +O\!\left(\log\frac1t\right).
\end{equation}
\end{corollary}
\begin{proof}
At \(r=\rho_j\), rewrite \eqref{vmb:eq:scale} as
\[
\frac{2e^{-x}}{t}=x+c+O(tx^2),\qquad x=jt.
\]
The leading equation has solution \(x=W(4/t)-c\), since
\(e^c=2\). Its derivative has order~\(x\), so the same two-stage
mean-value argument gives an error \(O(tx)\) in~\(x\), or \(O(x)\)
in~\(j\). Also \(x\asymp\log(1/t)\). Between consecutive real poles,
monotonicity sandwiches the counting function between the two adjacent
values. The leading expression in \eqref{vmb:eq:polecount} is decreasing
in~\(t\), so the same error estimate holds throughout the gaps.
\end{proof}

\section{Inverting the area-counting asymptotics}\label{vmb:sec:inverse}

\subsection{Finite-sector real interpolants}
Choose a cutoff \(R>\rho\) as in Theorem~\ref{vmb:thm:sectors}.
The finite principal-part expansion can be interpolated for real~\(x\)
as
\begin{equation}\label{vmb:eq:interpolant}
F_R(x)=Ce^{Lx}\bigl(1+E_R(x)\bigr),\qquad L=\log\lambda>0.
\end{equation}
For a nonreal conjugate pair, use its real oscillatory expression; for
a negative real pole use \(r^{-x}\cos(\pi x)\) in place of
\((-r)^{-n}\), with the same polynomial prefactor. This prescription
agrees with every integer coefficient contribution and is real analytic.
All nonleading terms and each fixed derivative are bounded by a
polynomial times \(e^{-\delta x}\), for some \(\delta>0\), because
there are only finitely many poles in the chosen expansion.
Consequently \(F_R\) is strictly increasing for sufficiently large~\(x\).
This is a specified asymptotic interpolant, not a claim that a discrete
sequence has a unique natural real interpolation.

\begin{theorem}[Inverse-sector expansion]\label{vmb:thm:inverse}
Let \(x_R(Y)\) solve \(F_R(x)=Y\) on its eventually increasing branch,
and put \(x_0=L^{-1}\log(Y/C)\). Set
\(\Phi(x)=-L^{-1}\log(1+E_R(x))\). Then
\begin{equation}\label{vmb:eq:Lagrange}
x_R(Y)=x_0+
\sum_{k\geq1}\frac1{k!}
\left(\frac{\dd}{\dd x}\right)^{k-1}
\Phi(x)^k\bigg|_{x=x_0}.
\end{equation}
For each fixed cutoff and sufficiently large~\(Y\), this is a
convergent small-correction expansion. In particular,
\begin{equation}\label{vmb:eq:inverse2}
x_R(Y)=x_0-\frac{E_R(x_0)}L
 +\frac{E_R(x_0)^2}{2L}
 +\frac{E_R(x_0)E_R'(x_0)}{L^2}
 +O\!\left(P(x_0)e^{-3\delta x_0}\right)
\end{equation}
for a suitable fixed polynomial~\(P\).
\end{theorem}
\begin{proof}
Writing \(x=x_0+h\) and taking logarithms in \eqref{vmb:eq:interpolant}
gives the exact equation \(h=\Phi(x_0+h)\). On a fixed small complex
neighborhood of~\(h=0\), \(\Phi\) and its derivative tend uniformly to
zero as \(x_0\to\infty\). Analytic Lagrange inversion applies; it may
also be obtained by inserting a parameter \(s\) into
\(h=s\Phi(x_0+h)\), expanding at~\(s=0\), and using the uniform smallness
to continue the series past \(s=1\). This proves \eqref{vmb:eq:Lagrange}.
Expanding \(\Phi=-E_R/L+E_R^2/(2L)+O(E_R^3)\), the terms with
\(k=1,2\) give \eqref{vmb:eq:inverse2}; all remaining terms have at least
three exponentially small factors, with polynomial derivative bounds.
\end{proof}

For one simple real sector \(E_R(x)=b e^{-\Delta x}\), the first terms
are especially transparent:
\begin{equation}\label{vmb:eq:singleinverse}
x_R(Y)=x_0-\frac bL\left(\frac YC\right)^{-\Delta/L}
 +\left(\frac1{2L}-\frac\Delta{L^2}\right)b^2
   \left(\frac YC\right)^{-2\Delta/L}+\cdots.
\end{equation}
For the \(\rho_k\) sector, \(b=C_k/C\) and
\(\Delta=\log(\rho_k/\rho)\). With several poles, products create
mixed exponent sums. One must retain any other sector of comparable or
larger size before interpreting a displayed correction as the next term
of an ordered asymptotic expansion.

\subsection{Integer thresholds and the rounding obstruction}
\begin{proposition}\label{vmb:prop:threshold}
The sequence \(d(n)\) is strictly increasing for \(n\geq1\). Define
\(N(Y)=\min\{n\geq1:d(n)\geq Y\}\). For each fixed cutoff~\(R\),
there exists \(K_R>0\) such that for all sufficiently large~\(Y\),
\begin{equation}\label{vmb:eq:ceilbounds}
\left\lceil x_R(Y)-K_RY^{-\eta_R}\right\rceil
\leq N(Y)\leq
\left\lceil x_R(Y)+K_RY^{-\eta_R}\right\rceil,
\qquad
\eta_R=\frac{\log(R/\rho)}{\log\lambda}>0.
\end{equation}
\end{proposition}
\begin{proof}
Append a part~1 to any valid composition. No new valley is created,
since the new step is nonascending; this is an injection from area~\(n\)
to area~\(n+1\). The singleton \((n+1)\) is not in its image, proving
strict increase.

At integers, Theorem~\ref{vmb:thm:sectors} gives
\(d(n)=F_R(n)+O_R(R^{-n})\). Near \(x_R(Y)\),
\(F_R'(x)\asymp Y\), so the corresponding uncertainty in the crossing
location is bounded by a constant times
\(R^{-x_R}/Y=O(Y^{-\eta_R})\).
One may make this precise first within a fixed distance of the crossing,
using the mean value theorem; outside that distance the strict leading
exponential separates the two sides. All integer points below the left
shift have count below~\(Y\), and all those above the right shift have
count at least~\(Y\). Taking ceilings yields the bounds.
\end{proof}

A single unconditional identity \(N(Y)=\lceil x_R(Y)\rceil\) would be
unjustified when \(Y\) lies arbitrarily close to an actual integer
threshold. Formula~\eqref{vmb:eq:ceilbounds} records the missing rounding
information rather than hiding it in a real-variable error term.

\begin{writenote}[Added 2 October 2026, batch~75: an instance of repository results]
Proposition~\ref{vmb:prop:threshold} and the remark after it are the
staircase theorem \texttt{p0:thm:staircase} of the transseries-and-inversion
volume \cite{vmb-tai}: part~(1) gives $N(Y)=\lceil\nu(Y)\rceil$ for any
increasing interpolation $\nu$ of the sequence, and part~(2), the
separation condition with $\eta=K_RY^{-\eta_R}$, says that the two ceilings
in \eqref{vmb:eq:ceilbounds} coincide, and then equal $N(Y)$, unless
$x_R(Y)$ lies within $\eta$ of an integer. The generic arithmetic of those
two parts is formalized in Lean as \texttt{Fabius.staircase\_ceil} and
\texttt{Fabius.staircase\_separation} in
\nolinkurl{Analysis/FabiusFunction/Lean/FabiusFunction/StaircaseInversion.lean},
statements about an arbitrary monotone interpolation that give no formal
status to this proposition. The series \eqref{vmb:eq:Lagrange} of
Theorem~\ref{vmb:thm:inverse} is Lagrange's reversion formula (in formal
form, \texttt{p0:thm:LB} of the same volume). No novelty is claimed for
these methods.
\end{writenote}

\section{Reproducible computations and proof checks}\label{vmb:sec:checks}

\subsection{Independent exact enumeration}
The package includes two logically independent finite checks.
First, a triangular coefficient algorithm implements
\eqref{vmb:eq:Urec}--\eqref{vmb:eq:Drec} using integers and formal-series
long division. Initializing \(U_{N+1}=D_{N+1}=0\) gives exactly all
coefficients through degree~\(N\); omitted objects have area larger
than~\(N\). Second, exhaustive enumeration generates every composition
through area~17, collapses consecutive equal parts, extracts interior
local minima, and checks their order directly. It agrees with the
coefficient algorithm. Lower bounds on the parts equal to~2 and~3 are
also checked independently through area~13.

The nonempty \(U_1\) coefficients agree with the displayed A001523
values through \(n=20\); its zero-degree coefficient differs only
because \(A=1+U_1\). The exact computation also gives
\[
d(100)=233832578565556052089825216845.
\]
The program retains all coefficients through~200 in
\texttt{verification.json}. Finite checks support the implementation;
they are not substitutes for the asymptotic proofs.

\subsection{Convergence of the finite ratios}
Let \(\varepsilon_n=d(n)/(C\lambda^n)-1\). High-precision numerical
comparison gives:
\begin{center}
\begin{tabular}{@{}rll@{}}
\toprule
\(n\)&\(d(n)/d(n-1)\)&\(\varepsilon_n\)\\
\midrule
10&2.0000000000000000000&\(-6.7798174\cdot10^{-2}\)\\
20&1.9808570502327331233&\(-9.1891465\cdot10^{-3}\)\\
30&1.9769260944011981878&\(-1.0612096\cdot10^{-3}\)\\
50&1.9764273761888468482&\(-1.4002503\cdot10^{-5}\)\\
75&1.9764207075861918763&\(-6.1918466\cdot10^{-8}\)\\
100&1.9764206780767492582&\(-2.7346103\cdot10^{-10}\)\\
150&1.9764206779458382883&\(-5.3327913\cdot10^{-15}\)\\
200&1.9764206779458357354&\(-1.0399431\cdot10^{-19}\)\\
\bottomrule
\end{tabular}
\end{center}
The finite-area quantities reproduce both estimates in
\eqref{vmb:eq:published}:
\begin{align*}
\frac{d(100)}{d(99)}
 &=1.976420678076749258223519681232\ldots,\\
\frac{d(100)}{(d(100)/d(99))^{100}}
 &=0.603901514896651955795442836948\ldots.
\end{align*}
This identifies the source of the discrepancy exactly; it is not a
change in the combinatorial definition or in the area indexing.

\subsection{Files and execution}
\texttt{verify.py} requires only the Python standard library. It checks
exact enumeration, agreement with the supplied classical initial terms,
and the rigorous rational intervals. Run
\begin{verbatim}
python verify.py
\end{verbatim}
The separate \texttt{numerics.py} requires \texttt{mpmath} and produces
high-precision exploratory values for poles, residues, and limit-law
constants. Its output explicitly distinguishes those numerical values
from certified enclosures. A small optional \texttt{expansion.py}
uses \texttt{sympy} to regenerate the \(\alpha_k(c)\) coefficients
from \eqref{vmb:eq:implicit-eps}. The article is self-contained and can be
compiled with
\begin{verbatim}
pdflatex article.tex
pdflatex article.tex
\end{verbatim}
No source font files, downloaded papers, or repository files are bundled.

\begin{writenote}[Added 2 October 2026, batch~75]
In the repository the programs live in \texttt{code/} and their recorded
outputs in \texttt{data/}. Each program writes its output file into its own
directory (\texttt{verify.py} writes \texttt{verification.json},
\texttt{numerics.py} writes \texttt{numerical\_results.json} and imports
\texttt{verify}, \texttt{expansion.py} writes
\texttt{expansion\_coefficients.json}), so the report's README runs them on
a scratch copy. The delivered build script and PDF are not shipped; the
README gives the build command.
\end{writenote}

\section{Further research questions and topics}\label{vmb:sec:future}

\subsection{The nonreal pole spectrum}
Determine the smallest-modulus nonreal poles, and certify their
multiplicities and ordering relative to \(\rho_2,\rho_3,\ldots\).
An effective route is to combine the explicit tail bounds with
argument-principle or Rouch\'e certificates on rationally described
contours. This would turn the abstract exponential cutoff in
Theorem~\ref{vmb:thm:sectors} into a numerically explicit first several
sectors, and could yield sharp eventual error signs or ratio
monotonicity. These latter conclusions have \emph{not} been inferred
from the numerical table.

\subsection{Is the unit circle a natural boundary?}
The real poles prove non-D-finiteness and obstruct meromorphic extension
through~1. They do not prove a natural boundary at every point of the
unit circle. A proof would need to locate noncancelling singularities
approaching a dense set of boundary points, or otherwise preclude local
continuation there. The factors \(Q_j\) and the root-of-unity structure
of the unimodal products give a concrete starting point.

\subsection{The next sector of pole accumulation}
The estimate \(x_j-X_j=O((\log j)^2/j)\) is beyond all logarithmic
orders. Determine its leading coefficient and a full expansion in
powers of~\(1/j\), with coefficient functions of~\(w_j\).
The proof suggests a systematic calculation: retain the \(k=2\) term
in \eqref{vmb:eq:logPexp}, Euler--Maclaurin corrections to
\((e^t-1)^{-1}\), and a sharper weighted telescoping expansion in
Lemma~\ref{vmb:lem:productcomparison}. The resulting two-scale transseries
should improve the pole-counting error beyond
\(O(\log(1/t))\).

\subsection{Joint limits and extreme heights}
The renewal representation can mark the number of columns, the number
of ones, peak heights, and the maximum height inside a repeated unit.
A joint normal law for additive statistics should follow after proving
nondegeneracy of the appropriate covariance matrix. Maximum height is
not additive: determine whether its logarithmic scale exhibits a lattice
Gumbel law and quantify any periodic phase. No extreme-value law is
claimed in this report.

\begin{writenote}[Added 5 October 2026, batch~98: answered in part]
The extreme-height half of this question is answered in Part~II:
Theorem~\ref{vmb:cw:thm:gumbel} proves
$\Prob(\mathsf H_n\le h)=\exp(-\alpha n\rho^h)(1+O(\log n/n))$ for the
maximum part $\mathsf H_n$ in the window $n\rho^h\asymp1$, with the explicit
intensity $\alpha=\rho^2P_1(\rho)/M$; \eqref{vmb:cw:eq:latticephase} gives the
lattice laws along phase subsequences, so there is no single continuous
Gumbel limit; and Theorem~\ref{vmb:cw:thm:independent} shows that the maximum
is asymptotically independent of the Gaussian bottom-valley count. The
sentence ``No extreme-value law is claimed in this report'' describes
Part~I. The first half, a joint normal law for additive statistics
(columns, ones, peak heights) with a nondegenerate covariance matrix, is not
addressed by Part~II and stays open
(Subsection~\ref{vmb:cw:sec:q-write}, item~7).
\end{writenote}

\subsection{Moving lower bounds and a crossover problem}
Theorem~\ref{vmb:thm:leading} holds for fixed~\(j\), while
Theorem~\ref{vmb:thm:W} describes large~\(j\). Determine a uniform estimate
for \(d_j(n)\) when \(j=j(n)\) grows. Since the pole gap shrinks as
\(j\) increases, a one-pole approximation need not remain uniform.
Locating the scale at which multiple poles interact would connect the
fixed-level exponential regime with sparse-composition regimes.

\subsection{Formalization and effective algorithms}
A useful ProveIt contribution would formalize the finite combinatorial
decomposition first, including the plateau convention, then develop a
general locally finite pole-product lemma. The exact interval
certificate is suitable for a small independently checked arithmetic
kernel. A separate algorithmic question is whether the exact coefficients
can be obtained substantially faster than triangular long division,
using the rapidly increasing minimum degrees of the kernels and
fast truncated multiplication. Non-P-recursiveness rules out a fixed
polynomial recurrence, not fast coefficient algorithms.

%% ====================================================================
%% Batch-98 addition (manuscript 04, "Critical Windows and Extreme
%% Heights in Valley-Monotone Compositions"). All labels of Part II carry
%% the prefix vmb:cw: (critical windows).
%% ====================================================================
\clearpage
\phantomsection\part*{Part II.\quad Critical windows and extreme heights}
\addcontentsline{toc}{part}{Part II.\texorpdfstring{\quad}{ } Critical windows and extreme heights}
\label{vmb:cw:part}

\noindent\emph{Sections~1--\ref{vmb:sec:future} and
Appendices~A--B form the original
report (Part~I) of 1~October 2026, unchanged except for dated notes. Part~II
was added on 5~October 2026 in batch~98 of ProveIt's incoming-report intake,
from one manuscript written against the pinned state of Part~I. It answers
the extreme-height question of Part~I (Section~\ref{vmb:sec:future},
``Joint limits and extreme heights''), introduces a valley-weighted model
with a first-order transition, and develops its critical windows.
Section~\ref{vmb:cw:sec:front} records the provenance and fixes the notation
of the addition; its appendices follow Part~I's.}

\section{Provenance, scope and notation of Part II}
\label{vmb:cw:sec:front}

\begin{writenote}[Added 5 October 2026, batch~98]
This section was written in the batch-98 write, except
Subsection~\ref{vmb:cw:sec:scope}, which prints the manuscript's own first
section with the changes listed there.
\end{writenote}

\subsection{Provenance}
\label{vmb:cw:sec:provenance}
Part~II is manuscript~04 of batch~98, the archive
\nolinkurl{OEIS_Bargraph_Critical_Windows_and_Extremes.zip}, which arrived in
commit \texttt{2172df76a} and was placed in this report by commit
\texttt{9353a7171}. Its article is titled \emph{Critical Windows and Extreme
Heights in Valley-Monotone Compositions}, with the subtitle ``An OEIS
A001523-linked continuation of ProveIt: phase coexistence, Dirichlet laws,
and invertible exponential expansions'' (24~pages, dated October~4, 2026);
its author line reads ``Research manuscript prepared for Vladimir
Reshetnikov''. AI assistance is not stated in the manuscript. Its pinned
repository commit is \texttt{70ab31a1ec77852d411bc8ca1fc81d86d0c9985e}
(4~October 2026), and the article blob it read,
\texttt{f37593705acecfdc8dc158e9079ae36887e96609}, is the text of Part~I as
committed in the batch-75 write (\texttt{9891e1f5e}); Part~I had not changed
between that pin and this write, so nothing in the manuscript is stale with
respect to it. The manuscript's source notes are shipped as
\texttt{02-critical-windows-SOURCE\_NOTES.md}; its programs and recorded
outputs are shipped with the prefix \texttt{02-critical-windows-} in
\texttt{code/} and \texttt{data/} (Section~\ref{vmb:cw:sec:checks}).

The manuscript cites OEIS A001523 (accessed October~4, 2026), the 2024
bargraph paper~\cite{FRV}, whose bibliographic record and OEIS linkage it
checked but whose full publisher text \emph{was not accessible} in its
session, Flajolet and Sedgewick~\cite{FS} (through the author-maintained page
\url{https://algo.inria.fr/flajolet/Publications/books.html}), and Part~I at
the pinned commit, which it treats as a research draft, ``not \ldots\ an
independently reviewed publication''. In Part~II these references are printed
as the entries \cite{OEIS}, \cite{FRV}, \cite{FS} of Part~I's bibliography,
with the manuscript's access notes kept in this paragraph, and its citations
of Part~I are printed as cross-references.

\subsection{What Part II proves, and how it relates to Part I}
\label{vmb:cw:sec:claims}
Part~I proves the exponential law, the positive-pole hierarchy, the
bottom-valley central limit theorem and a Boltzmann terminal limit for the
uniform model. Part~II keeps the same objects and adds:
\begin{itemize}[leftmargin=*,itemsep=1pt]
\item the \emph{lattice-Gumbel law} of the maximum part,
 $\Prob(\mathsf H_n\le h)=\exp(-\alpha n\rho^h)(1+O(\log n/n))$ in the
 window $n\rho^h\asymp1$, its lattice phase, and the asymptotic independence
 of $\mathsf H_n$ from the Gaussian bottom-valley count
 (Theorems~\ref{vmb:cw:thm:gumbel} and~\ref{vmb:cw:thm:independent});
\item a weighted model with three regimes in the bottom-valley weight~$u$,
 the exact critical weight
 $u_c=\rho_2(1-\rho_2)(1-\rho_2^2)/(1+\rho_2^4)$, a uniform $1/n$ critical
 window and its all-orders expansion
 (Theorems~\ref{vmb:cw:thm:three}, \ref{vmb:cw:thm:window},
 Proposition~\ref{vmb:cw:prop:windowall});
\item Dirichlet area laws at simultaneous critical weights on finitely many
 levels (Theorem~\ref{vmb:cw:thm:dirichlet}) and the Dirichlet-mixture law of
 the maximum part there (Theorem~\ref{vmb:cw:thm:mixed});
\item an exact $q$-product tail identity and the expansion of the
 height-cutoff pole to every exponential order
 (Lemma~\ref{vmb:cw:lem:psi}, Theorem~\ref{vmb:cw:thm:cutoffseries});
\item inverses of these expansions with rounding qualifications
 (Section~\ref{vmb:cw:sec:inversion}).
\end{itemize}
It also re-proves parts of Part~I so as to be independent of them.
Table~\ref{vmb:cw:tab:routes} lists every overlap. Where Part~II's statement
coincides with Part~I's, Part~I's statement stands and Part~II's text is a
second route; the generalizations are new.

\begin{table}[htbp]
\centering\small
\renewcommand{\arraystretch}{1.12}
\begin{tabular}{@{}>{\raggedright\arraybackslash}p{.30\textwidth}>{\raggedright\arraybackslash}p{.62\textwidth}@{}}
\toprule
Part II & Relation to Part I or to the repository\\
\midrule
Lemma~\ref{vmb:cw:lem:decomp}, \eqref{vmb:cw:eq:Usum} &
second route to Lemmas~\ref{vmb:lem:U}, \ref{vmb:lem:D} and
\eqref{vmb:eq:Usum} (the decomposition is originally~\cite{FRV}); the
weighted form \eqref{vmb:cw:eq:weightedrec} at every level is new and
generalizes \eqref{vmb:eq:marked}, which marks level~1 only\\
Lemma~\ref{vmb:cw:lem:poles} &
second route to Theorem~\ref{vmb:thm:meromorphic}, Lemma~\ref{vmb:lem:roots}
and Theorem~\ref{vmb:thm:leading} (same amplitudes $c_j$)\\
Theorem~\ref{vmb:cw:thm:three} &
its case $u=1$ (supercritical) is Theorem~\ref{vmb:thm:CLT}, and the
tightness of the final block there is part of
Theorem~\ref{vmb:thm:terminal}; the regimes $u\ne1$ are new\\
Proposition~\ref{vmb:cw:prop:rounding} &
its monotonicity argument is that of Proposition~\ref{vmb:prop:threshold};
its ceilings are an instance of \texttt{p0:thm:staircase}~\cite{vmb-tai}\\
\eqref{vmb:cw:eq:Wcrit}--\eqref{vmb:cw:eq:Wlog},
\eqref{vmb:cw:eq:Wmulti}--\eqref{vmb:cw:eq:multiinverseimplicit},
Theorem~\ref{vmb:cw:thm:inverseheight} &
instances of \texttt{p0:thm:lambert-core}, \texttt{p0:thm:core-reversion}
and \texttt{p0:thm:LB}~\cite{vmb-tai} (note in
Section~\ref{vmb:cw:sec:inversion}); the constants are new\\
everything else &
new in this report: Theorems~\ref{vmb:cw:thm:three} ($u\ne1$),
\ref{vmb:cw:thm:window}, \ref{vmb:cw:thm:dirichlet},
\ref{vmb:cw:thm:cutoffseries}, \ref{vmb:cw:thm:gumbel},
\ref{vmb:cw:thm:independent}, \ref{vmb:cw:thm:mixed},
\ref{vmb:cw:thm:inverseheight}, Proposition~\ref{vmb:cw:prop:windowall},
Lemmas~\ref{vmb:cw:lem:simplex}, \ref{vmb:cw:lem:psi},
Corollaries~\ref{vmb:cw:cor:firstorder}, \ref{vmb:cw:cor:cov}\\
\bottomrule
\end{tabular}
\caption{Overlaps of Part II with Part I and with the repository.}
\label{vmb:cw:tab:routes}
\end{table}

\paragraph{Answered and open questions of Part~I.}
Part~I's question ``Joint limits and extreme heights''
(Section~\ref{vmb:sec:future}) asked two things. Its second half, whether the
maximum height has a lattice Gumbel law and what its periodic phase is, is
answered by Theorem~\ref{vmb:cw:thm:gumbel} and~\eqref{vmb:cw:eq:latticephase},
and Theorem~\ref{vmb:cw:thm:independent} adds the joint law of the bottom-valley
count and the maximum. Its first half, a joint normal law for additive
statistics (columns, ones, peak heights) with a nondegenerate covariance
matrix, is \emph{not} addressed by Part~II and stays open. Part~I's questions
on the natural boundary and on moving lower bounds are restated, not
answered, in Section~\ref{vmb:cw:sec:future}.

\paragraph{What the write proved.}
Five statements of the manuscript are supported only by a sketch, an
analogy or a one-line remark. The write supplies proofs of them, each marked
\textsf{[write]} and printed with its proof: the tilted multicritical window
(Proposition~\ref{vmb:cw:prop:tilted}), the fixed-parameter phase-chamber
classification (Proposition~\ref{vmb:cw:prop:chambers}), the conditional
intensity of the critical maximum and its dependence on the area fractions
(Proposition~\ref{vmb:cw:prop:conditional}), the multicritical rounding and
critical quantile bounds (Proposition~\ref{vmb:cw:prop:roundingmulti}), and the
monotonicity of the interpolated cutoff shift
(Lemma~\ref{vmb:cw:lem:deltamono}). These proofs have not been independently
reviewed. Section~\ref{vmb:cw:sec:future} lists them with the manuscript's
open questions.

\subsection{Notation}
\label{vmb:cw:sec:notation}
Part~II uses Part~I's names for every shared object. The manuscript uses
several letters that Part~I uses for other objects, and overloads some
internally ($a$, $b$, $c$, $C$, $D$, $T$, $p$, $x_0$, $a_r$).
Table~\ref{vmb:cw:tab:notation} fixes the convention; every renaming is
listed, and no normalization is changed. The macros \verb|\Pp| and
\verb|\coef| of the manuscript are printed with Part~I's \verb|\Prob| and
\verb|\coeff|; the manuscript's unused macros \verb|\cO| and \verb|\tr| are
not carried over. Short identifications with Part~I set as
\wrt{\ldots} were added in writing on 5~October 2026, like the
\textsf{[write]} notes.

\begin{longtable}{@{}>{\raggedright\arraybackslash}p{.30\textwidth}>{\raggedright\arraybackslash}p{.24\textwidth}>{\raggedright\arraybackslash}p{.38\textwidth}@{}}
\caption{Notation of Part~II against the manuscript and Part~I.}
\label{vmb:cw:tab:notation}\\
\toprule
Notion & Part II & Manuscript; remark\\
\midrule
\endfirsthead
\toprule
Notion & Part II & Manuscript; remark (continued)\\
\midrule
\endhead
second positive pole & $\rho_2$ & $R$ (Part~I's $R$, $R_j$ are radii)\\
pole at the $m$-level critical point & $\rho_{m+1}$ & $R_m$; also ``$R$'' in its Section~8\\
radius beyond the dominant pole & $R$, $R_j$ (as in Theorem~\ref{vmb:thm:leading}) & $T$, $T_j$ (Part~I's $T$ is the limiting terminal block)\\
moments of $Z$ at $\rho=\rho_1$ & $M=M_\rho$, $V=V_\rho$ of \eqref{vmb:eq:MV} & $M_\rho$, $V_\rho$\\
critical moments at $\rho_2$ & $M_c$, $V_c$, $J_c$, $D_c=V_c/(2M_c^3)$ & $M$, $V$, $J$, $D$ (Part~I's $M,V$ are at $\rho$; Part~I's $D_j$ are series)\\
window amplitude & $d_c$ & $d$ (not Part~I's $d(n)$)\\
first window correction & $G_1(s)$ & $C_0(s)$ (equal by Proposition~\ref{vmb:cw:prop:windowall})\\
leading coefficient at $m$ levels & $\kappa^{(m)}$, with $\kappa^{(1)}=\kappa$ & $C_m$ (Part~I's $C_k$ are amplitudes at $\rho_k$)\\
critical counting polynomial & $\Pi_m(x)=\kappa^{(m)}x^m(1+\pi_1/x+\cdots)$ & $P_m$, $p_k$\\
area in level-$j$ repeated units & $\mathcal A_{j,n}$ & $C_{j,n}$\\
area of the final block & $|T_n|$ (Part~I's block $T_n$) & $T_n$\\
Dirichlet coordinates & $\xi_0,\ldots,\xi_m$, $\bm\xi$ & $x_0,\ldots,x_m$, $\bm x$ (Part~I's $x_j=jt_j$; the manuscript's $x_0$ is also a leading inverse)\\
area marks & $v_j$ & $t_j$ (Part~I's $t_j=-\log\rho_j$)\\
product $\prod_{r\ge j}(1-q^r)^{-2}$ & $P_j(q)$, $P_1(q)$ (Section~\ref{vmb:sec:accumulation}) & $\mathcal P_j$, $\mathcal P$ (Part~I's $\calP_R$ is a set of poles)\\
cutoff polynomials & $\varpi_{r-1}(h)$ & $Q_{r-1}(h)$ (Part~I's $Q_j$ is a series)\\
order-$r$ cutoff coefficients & $A_r(q)$ & $a_r(q)$ in a proof, $A_r(q)$ in Appendix~B (the $a_j(q)$ of \eqref{vmb:cw:eq:generaltail} are level-$j$ coefficients; $a_1=A_1$ only)\\
product coefficients & $\gamma_\ell(q)$ & $c_\ell(q)$ (Part~I's $c_j$ are amplitudes)\\
capped amplitude & $\hat c_h$ & $c_h$\\
level-2 numerator, Laurent constant & $B_2=E_2D_3$, $c_{2,0}$ & $A_2$, $d_{20}$\\
nearby-pole numerator & $\Xi_n$, $\Xi_n(\rho_{m+1})\to a_*$ & $A_n$, $a$\\
subcritical ratio & $\bar u=u/u_c$ & $a$\\
height-window bounds & $\tau_-\le n\rho^h\le\tau_+$ & $a\le n\rho^h\le b$\\
inverse shift & $a_c=\kappa_0/\kappa$ & $a$\\
height centres & $\beta_n$, $\beta_n^{\mathrm{crit}}$, $\beta_p(n)$ & $b_n$, $b_n^{\mathrm{crit}}$, $b_p(n)$ (not the Laurent $b$, $b_0$)\\
standard normal limit & $\chi$ & $Z$ (Part~I's $Z$ is the unit size)\\
multicritical inverse correction & $\varphi$ & $h$ (a height elsewhere)\\
number of minimal roots & $\nu$ & $p$\\
log-series variables in \eqref{vmb:cw:eq:Wlog} & $z_c$, $\ell_1$, $\ell_2$ & $Z$, $L_1$, $L_2$ (its $L_m$ at $m=1$ is $L_c$)\\
rounding constant & $C_*$ & $C$ (Part~I's $C$ is the growth amplitude)\\
\midrule
\multicolumn{3}{@{}l}{\emph{Kept, with a warning}}\\
marked function & $F(q,u)$ & $=D(q,u)$ of \eqref{vmb:eq:marked}; not the interpolant $F_R$\\
$\log(1/\rho)$ & $L$ & $=\log\lambda$, Part~I's $L$\\
Gumbel intensities & $\alpha$, $\alpha_j^*$ & not the coefficients $\alpha_k(c)$ of \eqref{vmb:eq:alphas}\\
weighted denominator & $Q_{j,u}$ & $Q_{j,1}=(1-q^j)Q_j$ with Part~I's $Q_j$\\
lattice phase & $\theta$ & not the ratio $\theta$ of \eqref{vmb:eq:main-intro}\\
cutoff shift & $\delta_h$, $\delta(h)$ & not the decay rate $\delta$ of \eqref{vmb:eq:interpolant}\\
run kernel, $(1-q^j)^{-2}$ & $H_j$, $E_j$ & $H_j$ as in Part~I; $\mathsf H_n$ is the maximum part; $E_j$ is not $E_R$\\
simplex & $\Delta_m$ & not the gap $\Delta$ of \eqref{vmb:eq:singleinverse}\\
dominant amplitude of $D_j$ & $c_j$ & as in Theorem~\ref{vmb:thm:leading}\\
\bottomrule
\end{longtable}

\subsection{The manuscript's research question, sources and contribution boundary}
\label{vmb:cw:sec:scope}
\begin{writenote}[Added 5 October 2026, batch~98]
This subsection is the manuscript's first section; its three subsections
are printed as unnumbered headings. Its citations of Part~I are printed as
cross-references, and nothing else is changed.
\end{writenote}

\subsubsection*{The OEIS connection is structural}
OEIS A001523 counts weakly unimodal compositions, or stacks. Its classical
asymptotics are already known, so the present article does not present them
as an open problem. Write
\begin{equation}
A(q)=\sum_{n\ge0}\mathrm{A001523}(n)q^n.
\end{equation}
The entry gives, among other formulas,
\begin{equation}\label{vmb:cw:eq:oeis}
 A(q)=1+\sum_{h\ge1}\frac{q^h}{1-q^h}
              \prod_{r=1}^{h-1}(1-q^r)^{-2}.
\end{equation}
It also links the 2024 paper of Fl\'orez, Ram\'irez, and Villamizar on
non-decreasing bargraphs \cite{OEIS,FRV}. The term ``non-decreasing'' refers
to the heights of successive interior valley plateaux, not to all columns.
We use the more explicit description \emph{valley-monotone}.

The repository report \emph{A Pole Hierarchy for Valley-Monotone Compositions}
(Part~I) rederives the minimum-part decomposition and proves its first
counting and probability consequences. Its final research section asks
whether maximum height has a lattice-Gumbel law and how joint statistics
behave. Those are specific, documented questions addressed here. The target
bargraph counting sequence is not identified with A001523: A001523 supplies
the unimodal blocks. No unverified OEIS number for the target sequence is
assigned.

\subsubsection*{What is inherited and what is developed here}
The symbolic decomposition, the unweighted dominant pole, the positive-pole
hierarchy, and the bottom-valley central limit theorem are inherited from
Part~I; the decomposition originates in \cite{FRV}. We give the short
proofs needed to make the present article independent of unproved repository
assertions. Meromorphic coefficient extraction, quasi-powers arguments,
Lagrange inversion, and the appearance of simplex laws from products of
simple poles are standard methods of analytic combinatorics \cite{FS}.

The present continuation develops the following model-specific conclusions:
the exact maximum-height intensity and its lattice phase; a weighted
first-order transition with a uniform critical density; its finite-size
scaling function and correction; a finite hierarchy of simultaneous critical
weights with Dirichlet area sharing; the resulting mixed extreme-height
law; and an explicit all-orders expansion and inversion for the moving
height-cutoff pole. The maximum-height theorem answers the explicit question
in Part~I. The weighted and multicritical formulations introduce and
solve further questions rather than claim to settle a previously advertised
OEIS conjecture.

\subsubsection*{Audit and limitations}
The repository source inspected was
\begin{quote}\small\ttfamily
SetTheory/Cardinals/docs/reports/\\
enumerative-combinatorics/valley-monotone-bargraphs/\\
article.tex
\end{quote}
at the snapshot identified by commit
\texttt{70ab31a1ec77852d411bc8ca1fc81d86d0c9985e}; the fetched article blob
was \texttt{f37593705acecfdc8dc158e9079ae36887e96609}.
The relevant definitions, analytic results, and final research questions
were read, not inferred from filenames. Targeted searches for the published
bargraph title and for maximum-height, phase-transition, and Dirichlet
refinements did not identify an earlier matching result. The publisher's
complete article was not accessible in this session; its bibliographic
record and the OEIS linkage were accessible. Consequently this is a bounded
source audit, not an exhaustive novelty certificate. All decimals introduced
below are numerical diagnostics unless explicitly described otherwise.

\section{Minimum-part decomposition and analytic foundations}
\label{vmb:cw:sec:foundations}

\subsection{Objects and generating functions}
A composition is a nonempty finite sequence of positive integers. Its area
is the sum of its parts, and its height is its largest part. Collapse every
maximal constant run to a single entry. An interior entry smaller than both
neighbors represents one valley plateau. A composition is valley-monotone
when the resulting valley heights are weakly increasing.

Let $\calU_j$ and $\calD_j$ be, respectively, the nonempty unimodal and
valley-monotone compositions with all parts at least $j$. Their ordinary
generating functions are $U_j(q)$ and $D_j(q)$. Thus $U_1=A-1$.
We use
\begin{equation}\label{vmb:cw:eq:kernels}
 H_j(q)=\frac{q^j}{1-q^j},\qquad
 E_j(q)=(1-q^j)^{-2},\qquad
 K_j(q)=H_j(q)U_{j+1}(q).
\end{equation}
The symbol $H_j$ counts a nonempty run at level $j$; it is not the random
maximum height, which will be denoted $\mathsf H_n$.

\begin{lemma}[Unique decomposition]\label{vmb:cw:lem:decomp}
As formal power series,
\begin{align}
 U_j&=H_j+E_jU_{j+1},\label{vmb:cw:eq:Urec}\\
 D_j&=H_j+\frac{E_jD_{j+1}}{1-K_j}.\label{vmb:cw:eq:Drec}
\end{align}
If each valley of height $j$ receives weight $u_j$, the corresponding
weighted function satisfies
\begin{equation}\label{vmb:cw:eq:weightedrec}
 D_j(q;u_j,u_{j+1},\ldots)
 =H_j(q)+\frac{E_j(q)D_{j+1}(q;u_{j+1},\ldots)}{1-u_jK_j(q)}.
\end{equation}
Only finitely many $u_j$ need differ from one.
\end{lemma}
\begin{proof}
A unimodal object either is an all-$j$ run or has a nonempty central object
with larger parts, flanked by two possibly empty $j$-runs. This proves
\eqref{vmb:cw:eq:Urec}.

For a valley-monotone object that is not all $j$, split at the runs of $j$.
Each nonfinal block of larger parts must be unimodal: an interior valley
there would have height larger than $j$ and would precede a valley of height
$j$. Conversely, choosing unimodal nonfinal blocks and a valley-monotone
final block creates no forbidden descent among valley heights. The unique
form is
\[
 (\text{initial }j\text{-run})\,
 (\calU_{j+1}\,\text{nonempty }j\text{-run})^*\,
 \calD_{j+1}\,(\text{final }j\text{-run}).
\]
Every repeated block contributes exactly one height-$j$ valley. Its kernel
is $K_j$, or $u_jK_j$ after marking. This proves both remaining formulas.
\end{proof}

\begin{writenote}[Added 5 October 2026, batch~98]
\eqref{vmb:cw:eq:Urec} and \eqref{vmb:cw:eq:Drec} are \eqref{vmb:eq:Urec} and
\eqref{vmb:eq:Drec} of Lemmas~\ref{vmb:lem:U} and~\ref{vmb:lem:D}, and this
proof is a second route to them. The weighted form
\eqref{vmb:cw:eq:weightedrec} at every level $j$ is new; at $j=1$ with
$u_2=u_3=\cdots=1$ it is \eqref{vmb:eq:marked}.
\end{writenote}

Iterating \eqref{vmb:cw:eq:Urec} gives
\begin{equation}\label{vmb:cw:eq:Usum}
 U_j(q)=\sum_{h\ge j}\frac{q^h}{1-q^h}
                   \prod_{r=j}^{h-1}(1-q^r)^{-2}.
\end{equation}
The unweighted target sequence $d(n)=\coeff qn D_1(q)$, with $d(0)=0$, begins
\[
1,2,4,8,16,32,64,128,256,512,1023,2042,4071,8106,16121,\ldots.
\]
The equality with $2^{n-1}$ through area ten is only an initial coincidence.

\subsection{A local analytic foundation, with no global conjectures}
\begin{lemma}[Dominant poles]\label{vmb:cw:lem:poles}
For every fixed $j$, $U_j$ is analytic on $|q|<1$. There is a unique
$\rho_j\in(0,1)$ such that $K_j(\rho_j)=1$, and
\[
 \rho_j<\rho_{j+1},\qquad \rho_j\longrightarrow1.
\]
The function $D_j$ is meromorphic on $|q|<1$. Its unique singularity on
$|q|\le\rho_j$ is a simple pole at $\rho_j$. For some $R_j>\rho_j$,
\begin{equation}\label{vmb:cw:eq:basepole}
 \coeff qn D_j=c_j\rho_j^{-n}+O(R_j^{-n}),\qquad
 c_j=\frac{E_j(\rho_j)D_{j+1}(\rho_j)}{M_j(\rho_j)}>0,
\end{equation}
where
\begin{equation}\label{vmb:cw:eq:moments}
 M_j(q)=\frac{qK_j'(q)}{K_j(q)},\qquad
 V_j(q)=q\frac{\dd}{\dd q}M_j(q).
\end{equation}
\end{lemma}
\begin{proof}
For $|q|\le r<1$, the product in \eqref{vmb:cw:eq:Usum} is uniformly bounded by
$\prod_{k\ge j}(1-r^k)^{-2}<\infty$, and the remaining summand is $O(r^h)$.
This gives local normal convergence. For real $q$, $K_j(q)$ increases
strictly from zero to infinity. Both $H_j>H_{j+1}$ and $U_{j+1}>U_{j+2}$,
so $K_j>K_{j+1}$ and the roots are strictly ordered.

For large $j$ and fixed $r<1$, all compositions with allowed parts at least
$j$ are bounded coefficientwise by $S_j/(1-S_j)$, where
$S_j(q)=q^j/(1-q)$ and $S_j(r)<1$. Thus $K_j=O(r^{2j})$ locally as
$j\to\infty$, which proves $\rho_j\to1$. Iteration of \eqref{vmb:cw:eq:Drec}
gives a locally normally convergent sum away from the locally finite zeros
of $1-K_j$: its multiplier is $E_j/(1-K_j)=1+O(r^j)$ and its additive
term is $H_j=O(r^j)$. This proves meromorphic continuation.

For $k>j$, $K_k(\rho_j)<1$, so the higher-level factors are analytic on
$|q|\le\rho_j$. The support of $K_j$ contains the consecutive degrees
$2j+1$ and $2j+2$. Equality in the triangle inequality in
$K_j(q)=1$, with $|q|=\rho_j$, therefore forces $q=\rho_j$. The pole is
simple because $K_j'(\rho_j)>0$, and its numerator is positive. Subtract
$c_j/(1-q/\rho_j)$ and apply Cauchy's coefficient bound on a slightly
larger circle containing no other singularity.
\end{proof}

\begin{writenote}[Added 5 October 2026, batch~98]
Lemma~\ref{vmb:cw:lem:poles} restates Theorem~\ref{vmb:thm:meromorphic},
Lemma~\ref{vmb:lem:roots} and Theorem~\ref{vmb:thm:leading}, and its proof is
a condensed second route to them. Its $\rho_j$ and $c_j$ are Part~I's: since
$K_j(\rho_j)=1$, $M_j(\rho_j)=\rho_jK_j'(\rho_j)$ and the two formulas for $c_j$
agree. The $S_j$ in the proof is Part~I's $S_J$ of
Proposition~\ref{vmb:prop:tail} with $J=j-1$.
\end{writenote}

For real $q\in(0,1)$, $M_j(q)$ and $V_j(q)$ are the mean and variance of
the positive integer distribution
\begin{equation}\label{vmb:cw:eq:boltzmannkernel}
 \Prob(Z_j=s)=\frac{\coeff qs K_j(q)\,q^s}{K_j(q)}.
\end{equation}
In particular $V_j(q)>0$. No natural-boundary assertion or classification of
nonreal poles is required in this paper. \wrt{at $j=1$, $q=\rho$ this is
Part~I's $Z$ of \eqref{vmb:eq:Z}, and $M_1(\rho)=M$, $V_1(\rho)=V$ are Part~I's
moments \eqref{vmb:eq:MV}.}

\section{A bottom-valley phase transition}
\label{vmb:cw:sec:transition}
Let $V_n$ count the valleys of height one. Give an area-$n$ object weight
$u^{V_n}$, with $u>0$, and normalize these weights to obtain a probability
measure $\Prob_{n,u}$. Define
\begin{equation}\label{vmb:cw:eq:F}
 F(q,u)=H_1(q)+\frac{B(q)}{1-uK(q)},\qquad
 K=K_1,\quad B=E_1D_2,\quad f_n(u)=\coeff qn F(q,u).
\end{equation}
\wrt{this $F$ is Part~I's $D(q,u)$ of \eqref{vmb:eq:marked}, with the same $B$
and $K$.} Put
\begin{equation}\label{vmb:cw:eq:critical}
 u_c=\frac1{K(\rho_2)},\quad
 M_c=M_1(\rho_2),\quad V_c=V_1(\rho_2),\quad
 J_c=\frac{u_c\rho_2^2K''(\rho_2)}2.
\end{equation}
At $\rho_2$ expand
\begin{equation}\label{vmb:cw:eq:Blaurent}
 B(q)=\frac b{1-q/\rho_2}+b_0+O(1-q/\rho_2),\qquad b>0.
\end{equation}
All these constants are defined analytically, not by their displayed decimals.

\begin{theorem}[Three regimes]\label{vmb:cw:thm:three}
For fixed $u>0$ the following statements hold.

\smallskip\noindent
\textbf{Subcritical regime, $u<u_c$.} Set $\bar u=u/u_c$. There is
$R>\rho_2$ such that
\begin{equation}\label{vmb:cw:eq:subcrit}
 f_n(u)=\frac b{1-\bar u}\rho_2^{-n}+O(R^{-n}),\qquad
 V_n\ \Longrightarrow\ \operatorname{Geom}_0(1-\bar u),
\end{equation}
where the limiting probability at $k\ge0$ is $(1-\bar u)\bar u^k$. The final
$\calD_2$ block has area $n-O_{\Prob}(1)$.

\smallskip\noindent
\textbf{Supercritical regime, $u>u_c$.} Let $r=r(u)<\rho_2$ solve $uK(r)=1$.
For some $R>r$,
\begin{equation}\label{vmb:cw:eq:supercrit}
 f_n(u)=\frac{B(r)}{M_1(r)}r^{-n}+O(R^{-n}),
\end{equation}
\begin{equation}\label{vmb:cw:eq:superCLT}
 \frac{V_n-n/M_1(r)}{\sqrt{nV_1(r)/M_1(r)^3}}
 \ \Longrightarrow\ \mathcal N(0,1).
\end{equation}
The final block has tight area, with an exponentially decaying limiting tail.

\smallskip\noindent
\textbf{Critical regime, $u=u_c$.} There is $R>\rho_2$ such that
\begin{equation}\label{vmb:cw:eq:criticalcount}
 f_n(u_c)=\rho_2^{-n}(\kappa n+\kappa_0)+O(R^{-n}),
\end{equation}
where
\begin{equation}\label{vmb:cw:eq:kappas}
 \kappa=\frac b{M_c},\qquad
 \kappa_0=\frac b{M_c}+\frac{b_0}{M_c}+\frac{bJ_c}{M_c^2}.
\end{equation}
Moreover,
\begin{equation}\label{vmb:cw:eq:uniform}
 \frac{V_n}{n}\Longrightarrow\Unif[0,1/M_c],\quad
 \frac{\E V_n}{n}\to\frac1{2M_c},\quad
 \frac{\Var V_n}{n^2}\to\frac1{12M_c^2}.
\end{equation}
\end{theorem}

\begin{proof}
When $u<u_c$, the denominator in \eqref{vmb:cw:eq:F} does not vanish on
$|q|\le\rho_2$; only the pole of $B$ contributes. Uniformly for complex $v$
near one,
\[
 \frac{f_n(uv)}{f_n(u)}\longrightarrow\frac{1-\bar u}{1-\bar uv}.
\]
This proves the geometric law. Marking the area outside the final block
by an extra variable close to one leaves all exterior factors analytic at
$\rho_2$; the same residue ratio gives a finite exponential moment for that
area, uniformly for large $n$.

For $u>u_c$, positivity and the consecutive kernel degrees give a unique
simple pole $r<\rho_2$. Its coefficient is $B(r)/M_1(r)$. Analytically
perturbing $u$ to $ue^s$ gives
\[
 \log\E_{n,u}e^{sV_n}
 =n\log\frac{r(u)}{r(ue^s)}+O(1),
\]
with an analytic, uniformly controlled $O(1)$ term. Differentiating
$\log K(r(ue^s))=-\log u-s$ gives first and second derivatives
$1/M_1(r)$ and $V_1(r)/M_1(r)^3$ for the leading logarithmic moment
generating function. Substitution $s=it/\sqrt n$ proves the central limit
theorem. Since $B$ is analytic beyond $r$, the final block is exponentially
tight.

At criticality, put $t=1-q/\rho_2$. Then
\[
 1-u_cK(q)=M_ct-J_ct^2+O(t^3).
\]
Using \eqref{vmb:cw:eq:Blaurent}, the singular part of $F$ is
\[
 \frac{b/M_c}{t^2}+\frac{b_0/M_c+bJ_c/M_c^2}{t}.
\]
Coefficient extraction proves \eqref{vmb:cw:eq:criticalcount}. The distributional
limit is proved with its full scaling window in Theorem~\ref{vmb:cw:thm:window}
below. Since every valley uses at least one part, $0\le V_n/n\le1$;
weak convergence therefore also implies convergence of the first two
moments and the variance statement.
\end{proof}

\begin{writenote}[Added 5 October 2026, batch~98]
The ordinary model $u=1$ lies in the supercritical regime, where
$r(1)=\rho$ and \eqref{vmb:cw:eq:superCLT} is Theorem~\ref{vmb:thm:CLT} with
$\mu=1/M$, $\sigma^2=V/M^3$; the tightness of the final block at $u=1$ is
part of Theorem~\ref{vmb:thm:terminal}. The regimes $u\le u_c$ and the
supercritical regime at $u\ne1$ are new.
\end{writenote}

\begin{corollary}[A first-order transition]\label{vmb:cw:cor:firstorder}
The exponential growth rate is
\begin{equation}
 \Lambda(u)=\lim_{n\to\infty}\frac1n\log f_n(u)
 =\begin{cases}-\log \rho_2,&u\le u_c,\\-\log r(u),&u>u_c.
 \end{cases}
\end{equation}
Its derivative with respect to $\log u$ jumps from $0$ to $1/M_c$ at $u_c$.
\end{corollary}
\begin{proof}
The coefficient asymptotics give the formula. Implicit differentiation of
$uK(r(u))=1$ gives $-\dd\log r/\dd\log u=1/M_1(r)$; now let $u\downarrow u_c$.
\end{proof}

The critical value has a compact exact expression in the second unweighted
pole:
\begin{equation}\label{vmb:cw:eq:ucclosed}
 \boxed{\displaystyle
 u_c=\frac{\rho_2(1-\rho_2)(1-\rho_2^2)}{1+\rho_2^4}.}
\end{equation}
Indeed, $K_2(\rho_2)=1$ gives $U_3(\rho_2)=(1-\rho_2^2)/\rho_2^2$, and
\eqref{vmb:cw:eq:Urec} then gives
$U_2(\rho_2)=(1+\rho_2^4)/(\rho_2^2(1-\rho_2^2))$.
Numerically,
\begin{align*}
 \rho_2&=0.6285205937788026\ldots,&u_c&=0.1221810090081765\ldots,\\
 M_c&=13.41194879035466\ldots,&1/M_c&=0.07456038012307058\ldots,\\
 \kappa&=0.2326297632869357\ldots,&
 \kappa_0&=-0.3327584620888633\ldots.
\end{align*}
The ordinary uniform model is $u=1>u_c$; the new transition is a weighted
refinement, not a claim that the ordinary model is critical.

\section{The critical window and all finite-size orders}
\label{vmb:cw:sec:window}
Define the entire function
\begin{equation}\label{vmb:cw:eq:G}
 G(z)=\int_0^1e^{zx}\,\dd x
 =\begin{cases}(e^z-1)/z,&z\ne0,\\1,&z=0.\end{cases}
\end{equation}
The integral definition is preferable near $z=0$.

\begin{theorem}[Uniform critical window]\label{vmb:cw:thm:window}
Let $u_n(s)=u_c e^{s/n}$. Uniformly for $s$ in any fixed compact subset
of $\C$,
\begin{equation}\label{vmb:cw:eq:window}
 \rho_2^nf_n(u_n(s))=\kappa nG(s/M_c)+G_1(s)+O(n^{-1}),
\end{equation}
where
\begin{align}
 D_c&=\frac{V_c}{2M_c^3},\qquad
 d_c=b\left(-\frac12+\frac1{M_c}+\frac {J_c}{M_c^2}\right)+\frac{b_0}{M_c},
 \label{vmb:cw:eq:Dd}\\
 G_1(s)&=\frac b2+e^{s/M_c}(d_c+bD_cs).\label{vmb:cw:eq:C0}
\end{align}
For real $s$, under $\Prob_{n,u_n(s)}$ the random variable $V_n/n$ converges
to the density
\begin{equation}\label{vmb:cw:eq:tiltdensity}
 p_s(x)=\frac{M_c e^{sx}}{G(s/M_c)}\,
 \mathbf1_{[0,1/M_c]}(x).
\end{equation}
In particular $p_0$ is the uniform density in \eqref{vmb:cw:eq:uniform}.
\end{theorem}

\begin{proof}
For $u$ near $u_c$, choose a fixed circle $|q|=R>\rho_2$ inside which the only
singularities of $F(q,u)$ are the pole at $\rho_2$ and the nearby root $r(u)$ of
$uK(r)=1$. Their grouped contribution is analytic in $u$, including at
coalescence. When the roots are distinct, residue extraction gives
\begin{equation}\label{vmb:cw:eq:twopoles}
 \rho_2^nf_n(u_c e^\epsilon)
 =\frac{b}{1-e^\epsilon}
 +\frac{B(r)}{M_1(r)}\exp\!\left(n\log\frac{\rho_2}r\right)
 +O((\rho_2/R)^n).
\end{equation}
The error is uniform near $\epsilon=0$: subtract the grouped principal
parts, or equivalently use a fixed contour encircling both poles together.

Write $r=\rho_2(1-\delta)$. Solving
$u_cK(\rho_2(1-\delta))=e^{-\epsilon}$ gives
\begin{equation}\label{vmb:cw:eq:deltacrit}
 \delta=\frac\epsilon{M_c}+
 \left(\frac {J_c}{M_c^3}-\frac1{2M_c}\right)\epsilon^2+O(\epsilon^3).
\end{equation}
Consequently
\begin{equation}\label{vmb:cw:eq:logrcrit}
 \log\frac{\rho_2}r=\frac\epsilon{M_c}+D_c\epsilon^2+O(\epsilon^3),
 \qquad
 \frac{B(r)}{M_1(r)}=\frac b\epsilon+d_c+O(\epsilon).
\end{equation}
Here $V_c=2J_c+M_c-M_c^2$. Also
$b/(1-e^\epsilon)=-b/\epsilon+b/2+O(\epsilon)$.
Insert $\epsilon=s/n$ into \eqref{vmb:cw:eq:twopoles}; the singular terms cancel
and give \eqref{vmb:cw:eq:window}--\eqref{vmb:cw:eq:C0}. The cancellation is holomorphic,
so the expansion is uniform also at $s=0$ and at complex zeros of $G$.
It is an additive estimate, not a relative estimate at such zeros.

For real $s,t$, the moment generating function of $V_n/n$ is
\[
 \frac{f_n(u_c e^{(s+t)/n})}{f_n(u_c e^{s/n})}
 \longrightarrow\frac{G((s+t)/M_c)}{G(s/M_c)}.
\]
This is exactly the transform of \eqref{vmb:cw:eq:tiltdensity}. The bounded support
of $V_n/n$ makes transform convergence sufficient and gives moment
convergence as well.
\end{proof}

\begin{proposition}[All-orders window expansion]\label{vmb:cw:prop:windowall}
For every $N\ge1$ there are entire functions $G_0,\ldots,G_N$ such that,
uniformly on fixed compact $s$-sets,
\begin{equation}\label{vmb:cw:eq:allwindow}
 \rho_2^nf_n(u_c e^{s/n})
 =\sum_{j=0}^{N}n^{1-j}G_j(s)+O(n^{-N}).
\end{equation}
Here $G_0(s)=\kappa G(s/M_c)$, $G_1$ is \eqref{vmb:cw:eq:C0}, and $G_j(0)=0$ for $j\ge2$.
Every coefficient is obtained by finite series arithmetic from the local
Taylor coefficients of $K$ and the Laurent coefficients of $B$ at $\rho_2$.
\end{proposition}
\begin{proof}
Expand the analytic inverse root $r(u_c e^\epsilon)$, the Laurent series
$B(r)/M_1(r)$, and $b/(1-e^\epsilon)$ in \eqref{vmb:cw:eq:twopoles} to a prescribed
finite order. The exponential $\exp(n\log(\rho_2/r))$ is then expanded after
$\epsilon=s/n$. At each order only finitely many Taylor coefficients are
needed. The simple $1/\epsilon$ terms cancel; all remaining coefficient
functions are polynomials in $s$ times $e^{s/M_c}$, together with a possible
removable divided difference. They extend to entire functions. Analytic
Taylor remainder bounds on a fixed $\epsilon$-disk prove the stated uniform
remainder. At $s=0$, \eqref{vmb:cw:eq:criticalcount} has exactly a degree-one
polynomial prefactor plus an exponentially small error, forcing all later
coefficients to vanish.
\end{proof}

This is a genuine uniform coalescing-pole expansion. Treating the two
residues separately and ignoring their cancellation would give divergent
prefactors and an incorrect critical approximation.

\section{A hierarchy of multicritical Dirichlet laws}
\label{vmb:cw:sec:multi}
The one-weight transition is the first member of a finite-level hierarchy.
Fix $m\ge1$, mark valleys of heights $1,\ldots,m$, and leave all higher
valleys unweighted. Set
\begin{equation}\label{vmb:cw:eq:multiweights}
 u_j^*=\frac1{K_j(\rho_{m+1})},\qquad
 M_j^*=M_j(\rho_{m+1}),\qquad 1\le j\le m.
\end{equation}
Iteration of the weighted decomposition gives the exact identity
\begin{equation}\label{vmb:cw:eq:multidecomp}
 F_m(q,\bm u)=
 \sum_{\ell=1}^{m}H_\ell(q)
   \prod_{j=1}^{\ell-1}\frac{E_j(q)}{1-u_jK_j(q)}
 +D_{m+1}(q)\prod_{j=1}^{m}\frac{E_j(q)}{1-u_jK_j(q)}.
\end{equation}
The last term is the branch that reaches level $m+1$. It has $m+1$ critical
factors; the earlier branches have at most $m-1$.
\wrt{at $m=1$, $\rho_{m+1}=\rho_2$, $u_1^*=u_c$, $M_1^*=M_c$ and $F_1=F$ of
\eqref{vmb:cw:eq:F}.}

\subsection{The elementary simplex mechanism}
Write
\[
 \Delta_m=\{(\xi_1,\ldots,\xi_m):\xi_j\ge0,\ \textstyle\sum_{j=1}^m \xi_j\le1\},
 \qquad \xi_0=1-\sum_{j=1}^m \xi_j.
\]
Integration on $\Delta_m$ means ordinary $m$-dimensional Lebesgue measure;
its volume is $1/m!$.

\begin{lemma}[Nearby simple poles]\label{vmb:cw:lem:simplex}
Suppose $r_{i,n}=R_0\exp(\eta_i/n+O(\log n/n^2))$, $0\le i\le m$, where
the real $\eta_i$ range over a compact set. Let $\Xi_n$ be
analytic and uniformly bounded on a fixed disk $|q|\le R>R_0$, with
$\Xi_n(R_0)=a_*+O(\log n/n)$. Then
\begin{align}\label{vmb:cw:eq:simplexlemma}
 &\coeff qn\left[\Xi_n(q)\prod_{i=0}^{m}(1-q/r_{i,n})^{-1}\right]\notag\\
 &\hspace{8pt}=R_0^{-n}n^m\left[
 a_*\int_{\Delta_m}\exp\left(-\sum_{i=0}^{m}\eta_i \xi_i\right)\dd\bm \xi
 +O\left(\frac{\log n}{n}\right)\right].
\end{align}
If the root and amplitude errors omit $\log n$, so does the final error.
Coincident $\eta_i$ cause no exception.
\end{lemma}
\begin{proof}
First take $\Xi_n=1$. The coefficient is the finite sum over
$k_0+\cdots+k_m=n$, $k_i\ge0$, of $\prod r_{i,n}^{-k_i}$. Dividing by
$R_0^{-n}$ gives, uniformly for every summand,
\[
 \exp\left(-\sum_i\eta_i k_i/n\right)
 \left(1+O(\log n/n)\right).
\]
The grid sum divided by $n^m$ is a Riemann sum on $\Delta_m$. Its integrand
and first derivatives are uniformly bounded, so its error is $O(1/n)$.
For general $\Xi_n$, convolve with its Taylor coefficients. Cauchy's bound
makes their absolute values at radius $R_0$ geometrically summable, including
one extra factor of the index. Shifting $n$ by that index changes the leading
Riemann sum by $O((k+1)n^{m-1})$. Summation therefore multiplies the leading
term by $\Xi_n(R_0)$ and adds only $O(R_0^{-n}n^{m-1})$. The assumed error in
$\Xi_n(R_0)$ completes the proof.
\end{proof}

\begin{writenote}[Added 5 October 2026, batch~98]
The manuscript states this lemma with its letters $R$ for the common pole
location and $T$ for the radius, printed here as $R_0$ and $R$; it is applied
at $R_0=\rho_{m+1}$, and $R_0=\rho_2$ when $m=1$.
\end{writenote}

\begin{theorem}[Multicritical count and area sharing]\label{vmb:cw:thm:dirichlet}
At the weights \eqref{vmb:cw:eq:multiweights}, for some $R>\rho_{m+1}$,
\begin{equation}\label{vmb:cw:eq:multicount}
 \coeff qn F_m(q,\bm u^*)=\rho_{m+1}^{-n}\Pi_m(n)+O(R^{-n}),
\end{equation}
where $\Pi_m$ is a degree-$m$ polynomial with positive leading coefficient
\begin{equation}\label{vmb:cw:eq:multilead}
 \kappa^{(m)}=\frac{c_{m+1}}{m!}
       \prod_{j=1}^{m}\frac{E_j(\rho_{m+1})}{M_j^*}.
\end{equation}
Let $\mathcal A_{j,n}$ be the total area in repeated units at level $j$, let
$|T_n|$ be the area of the final $\calD_{m+1}$ block, and let $V_{j,n}$ count
height-$j$ valleys. On the dominant branch,
\begin{equation}\label{vmb:cw:eq:dirichlet}
 \left(\frac{\mathcal A_{1,n}}n,\ldots,\frac{\mathcal A_{m,n}}n,\frac{|T_n|}n\right)
 \Longrightarrow\Dir(1,\ldots,1),
\end{equation}
with $m+1$ coordinates, and
\begin{equation}\label{vmb:cw:eq:countsareas}
 \frac{M_j^*V_{j,n}-\mathcal A_{j,n}}n\longrightarrow0
 \quad\text{in probability}.
\end{equation}
The total endpoint-run area is tight. The probability of an earlier branch
is $O(n^{-2})$ (exponentially small when $m=1$).
\end{theorem}
\begin{proof}
At $\rho_{m+1}$, the tail $D_{m+1}$ has one simple pole and each denominator in
\eqref{vmb:cw:eq:multidecomp} has one simple zero. The highest Laurent coefficient
is
$c_{m+1}\prod E_j(\rho_{m+1})/M_j^*$, giving pole order $m+1$ and
\eqref{vmb:cw:eq:multilead}. The other branches have pole order at most $m-1$.
There are no other singularities on the first circle; subtraction of the
finite Laurent principal part proves \eqref{vmb:cw:eq:multicount}. This also proves
the earlier-branch probability estimate.

Mark the area in the level-$j$ repeated units by replacing $K_j(q)$ with
$K_j(qe^{v_j/n})$, and mark terminal area by replacing $D_{m+1}(q)$ with
$D_{m+1}(qe^{v_0/n})$. Their poles move to
$\rho_{m+1}e^{-v_j/n}$ and $\rho_{m+1}e^{-v_0/n}$. The remaining analytic factors are
uniformly bounded and converge to their original values. Applying
Lemma~\ref{vmb:cw:lem:simplex} and normalizing gives the moment generating function
\[
 m!\int_{\Delta_m}\exp\left(\sum_{i=0}^m v_i\xi_i\right)\dd\bm \xi,
\]
which is that of the stated Dirichlet distribution.

If valley counts are simultaneously marked by $e^{s_j/n}$, the $j$th pole
moves instead by logarithmic amount
$-(v_j+s_j/M_j^*)/n+O(n^{-2})$. The joint limiting transform therefore
identifies $V_{j,n}/n$ with $\xi_j/M_j^*$ and $\mathcal A_{j,n}/n$ with $\xi_j$.
This proves \eqref{vmb:cw:eq:countsareas}. Endpoint-run marking affects only
analytic factors; a small fixed positive exponential mark is allowed
because $\rho_{m+1}<1$. The corresponding normalized coefficient ratio stays
bounded, proving tightness with an exponential moment.
\end{proof}

\begin{writenote}[Added 5 October 2026, batch~98]
At $m=1$, $c_2E_1(\rho_2)=c_2/(1-\rho_2)^2=b$ by
Appendix~\ref{vmb:cw:app:laurent}, so $\kappa^{(1)}=b/M_c=\kappa$ and
$\Pi_1(n)=\kappa n+\kappa_0$, consistent with
\eqref{vmb:cw:eq:criticalcount}; for $m=1$ there is a single earlier branch,
$H_1$, which is analytic at $\rho_2$.
\end{writenote}

\begin{corollary}[Macroscopic covariance]\label{vmb:cw:cor:cov}
At the fully critical weights,
\begin{align}
 \frac{\E V_{j,n}}n&\longrightarrow\frac1{(m+1)M_j^*},\\
 \frac{\Var V_{j,n}}{n^2}&\longrightarrow
 \frac{m}{(m+1)^2(m+2)(M_j^*)^2},\\
 \frac{\Cov(V_{i,n},V_{j,n})}{n^2}&\longrightarrow
 -\frac1{(m+1)^2(m+2)M_i^*M_j^*}\quad(i\ne j).
\end{align}
\end{corollary}
\begin{proof}
The count densities are bounded. Their moments converge to those of
$\xi_j/M_j^*$ under the uniform-simplex distribution. Integrating monomials
on $\Delta_m$ gives the displayed means and covariances.
\end{proof}

\subsection{Finite-dimensional critical windows and phase chambers}
For $u_{j,n}=u_j^*e^{s_j/n}$, the normalized leading coefficient is multiplied
by
\begin{equation}\label{vmb:cw:eq:multiwindow}
 \mathcal G_m(\bm s)=m!\int_{\Delta_m}
 \exp\left(\sum_{j=1}^m\frac{s_j\xi_j}{M_j^*}\right)\dd\bm \xi.
\end{equation}
The limiting area law is the Dirichlet law tilted by this exponential.
The proof is the same nearby-pole argument. In particular the critical
window has width $1/n$ in every logarithmic weight coordinate.

\begin{proposition}[{[write]} The tilted multicritical window]\label{vmb:cw:prop:tilted}
Let $\bm s$ range over a compact subset of $\mathbb R^m$ and
$u_{j,n}=u_j^*e^{s_j/n}$. Then
\[
 \frac{\coeff qn F_m(q,\bm u_n)}{\coeff qn F_m(q,\bm u^*)}
 =\mathcal G_m(\bm s)+O(n^{-1}),
\]
and, under $\Prob_{n,\bm u_n}$, the vector in \eqref{vmb:cw:eq:dirichlet}
converges in law to the probability measure with density proportional to
$\exp(\sum_{j=1}^m s_j\xi_j/M_j^*)$ on $\Delta_m$ (with respect to
$\dd\bm\xi$), and \eqref{vmb:cw:eq:countsareas} holds.
\end{proposition}
\begin{proof}
This proof was added in writing; the manuscript gives the sentence before
the proposition. Each $u_j^*K_j$ has a simple zero of $1-u_j^*K_j$ at
$\rho_{m+1}$, with $\dd\log(u_j^*K_j)/\dd\log q=M_j^*>0$ there. By the
analytic implicit function theorem, for $\epsilon$ near~$0$ the equation
$u_j^*e^{\epsilon}K_j(r)=1$ has a unique root near $\rho_{m+1}$,
$\log(r/\rho_{m+1})=-\epsilon/M_j^*+O(\epsilon^2)$, and
$1-u_j^*e^{\epsilon}K_j(q)=(1-q/r)g_j(q,\epsilon)$ with $g_j$ analytic and
nonvanishing on a fixed disk $|q|\le R$, $R>\rho_{m+1}$, uniformly in
$\epsilon$ (choose $R$ so that no other zero of $1-u_j^*K_j$ and no other
pole of $D_{m+1}$ lies in $|q|\le R$; both sets are locally finite, and
$\rho_{m+1}$ is the only singular point on $|q|=\rho_{m+1}$ by the
triangle-inequality argument of Lemma~\ref{vmb:cw:lem:poles}). Write
$D_{m+1}(q)=(1-q/\rho_{m+1})^{-1}g_0(q)$ likewise. With
$\epsilon=s_j/n$, the dominant branch of \eqref{vmb:cw:eq:multidecomp} is
$\Xi_n(q)\prod_{i=0}^m(1-q/r_{i,n})^{-1}$ with $r_{0,n}=\rho_{m+1}$,
$\eta_0=0$, $\eta_j=-s_j/M_j^*$, root errors $O(n^{-2})$ and
$\Xi_n=g_0\prod_jE_j/g_j$, whose value at $\rho_{m+1}$ is
$a_*+O(1/n)$ with $a_*=c_{m+1}\prod_jE_j(\rho_{m+1})/M_j^*=m!\,\kappa^{(m)}$.
Lemma~\ref{vmb:cw:lem:simplex} (without logarithms) gives
$\rho_{m+1}^{-n}n^m[a_*\int_{\Delta_m}\exp(\sum_js_j\xi_j/M_j^*)\dd\bm\xi
+O(1/n)]$. The earlier branches have at most $m-1$ such factors, and the
same lemma bounds their coefficients by $O(\rho_{m+1}^{-n}n^{m-2})$. Dividing
by the case $\bm s=0$, whose integral is $1/m!$, gives the first claim.
Adding the area marks $v_j/n$, $v_0/n$ of the proof of
Theorem~\ref{vmb:cw:thm:dirichlet} shifts $\eta_i$ by $-v_i$; the ratio of the
marked to the unmarked coefficient converges to
$\int\exp(\sum_iv_i\xi_i+\sum_js_j\xi_j/M_j^*)\dd\bm\xi/
\int\exp(\sum_js_j\xi_j/M_j^*)\dd\bm\xi$ for all real $\bm v$, which is the
moment generating function of the tilted law; the variables are bounded,
so this identifies the limit. The count-area identification is that of
Theorem~\ref{vmb:cw:thm:dirichlet}, with $\eta_j$ shifted by the constant
$-s_j/M_j^*$.
\end{proof}

There is also a complete fixed-parameter classification. For arbitrary
positive weights $u_1,\ldots,u_m$, let $r_j$ solve $u_jK_j(r_j)=1$ and set
\begin{equation}\label{vmb:cw:eq:chambers}
 r_* =\min(r_1,\ldots,r_m,\rho_{m+1}).
\end{equation}
If precisely $\nu$ of these $m+1$ values equal $r_*$, then the dominant pole
has order $\nu$ and the coefficients have form
$r_*^{-n}$ times a degree-$(\nu-1)$ polynomial, plus an exponentially smaller
error. To see noncancellation, take the suffix after the last active
weighted level: all its inactive denominators are positive and finite at
$r_*$. Its positive value, or its positive tail residue, is multiplied by
exactly the active simple-pole factors. Earlier branches cannot remove this
highest-order positive contribution. Thus the phase boundaries are the
explicit equal-root hypersurfaces, not hypothetical singularities inferred
from a numerical phase plot.

\begin{proposition}[{[write]} Phase chambers]\label{vmb:cw:prop:chambers}
In the setting of \eqref{vmb:cw:eq:chambers}, $r_*$ is the only singularity
of $F_m(q,\bm u)$ on $|q|\le r_*$; it is a pole of order exactly $\nu$, and
for some $R>r_*$,
$\coeff qn F_m(q,\bm u)=r_*^{-n}\Pi(n)+O(R^{-n})$ with a polynomial $\Pi$ of
degree $\nu-1$ and positive leading coefficient.
\end{proposition}
\begin{proof}
This proof was added in writing, expanding the manuscript's sketch above.
Each $r_j\in(0,1)$ exists and is unique, because $K_j$ increases from $0$ to
$\infty$ on $(0,1)$. Call an index active if its value in
\eqref{vmb:cw:eq:chambers} equals $r_*$ (the index $m+1$ standing for
$\rho_{m+1}$). For $|q|\le r_*$ and $j\le m$,
$|u_jK_j(q)|\le u_jK_j(|q|)\le u_jK_j(r_j)=1$, with equality only if
$|q|=r_j=r_*$, and then, by the consecutive degrees $2j+1,2j+2$ in the
support of $K_j$ (Lemma~\ref{vmb:cw:lem:poles}), only at $q=r_*$. So $1-u_jK_j$
has no zero on $|q|\le r_*$ for inactive $j$ and only the simple zero $r_*$
for active $j$ (simple because $K_j'(r_*)>0$). By
Lemma~\ref{vmb:cw:lem:poles}, $D_{m+1}$ is analytic on $|q|\le r_*$ unless
$m+1$ is active, and then its only singularity there is the simple pole
$\rho_{m+1}=r_*$ with residue amplitude $c_{m+1}>0$. The functions $H_\ell$,
$E_j$ are analytic on the unit disk. This proves the first claim, and the
order of the pole of each branch of \eqref{vmb:cw:eq:multidecomp} is the
number of active indices among its factors: $\#(\mathcal I\cap\{1,\ldots,\ell-1\})$
for the branch $\ell\le m$, and $\nu$ for the final branch, where
$\mathcal I$ is the active set. So the order of $F_m$ at $r_*$ is at most
$\nu$, and the coefficient of $(1-q/r_*)^{-\nu}$ is the sum, over the branches
attaining order $\nu$, of products of: $H_\ell(r_*)>0$ or the value
$D_{m+1}(r_*)>0$ (inactive $m+1$) or the amplitude $c_{m+1}>0$ (active
$m+1$); $E_j(r_*)>0$; $1/(1-u_jK_j(r_*))>0$ for inactive $j$ (since
$u_jK_j(r_*)<u_jK_j(r_j)=1$); and $1/M_j(r_*)>0$ for active $j$, the
leading coefficient of $1/(1-u_jK_j(q))$ at $r_*$ because
$u_jr_*K_j'(r_*)=M_j(r_*)$ when $u_jK_j(r_*)=1$. The final branch always
attains order $\nu$, so this sum is positive, the order is exactly $\nu$,
and no cancellation occurs. Poles of the factors are locally finite, so a
radius $R>r_*$ with no other singularity on $|q|\le R$ exists; subtracting
the principal part at $r_*$ and applying Cauchy's bound on $|q|=R$ gives the
expansion, with $\coeff qn(1-q/r_*)^{-k}=\binom{n+k-1}{k-1}r_*^{-n}$ giving
the polynomial of degree $\nu-1$.
\end{proof}

\wrt{at $m=1$ the chambers are those of Theorem~\ref{vmb:cw:thm:three}:
$\nu=1$ off $u_1=u_c$ and $\nu=2$ at $u_1=u_c$, where $r_1=\rho_2$.}

\section{An exact height-tail identity and its exponential expansion}
\label{vmb:cw:sec:heighttail}
Let $U_j^{\le h}$ and $D_j^{\le h}$ count the same objects with every part
at most $h$. For $h\ge j$,
\begin{equation}\label{vmb:cw:eq:Utail}
 U_j-U_j^{\le h}=U_{h+1}\prod_{r=j}^{h}(1-q^r)^{-2}.
\end{equation}
This is simply the tail of the maximum-part decomposition.
Put
\begin{equation}\label{vmb:cw:eq:Pj}
 P_j(q)=\prod_{r=j}^{\infty}(1-q^r)^{-2},\qquad
 P_1(q)=\prod_{r=1}^{\infty}(1-q^r)^{-2}.
\end{equation}
These products are analytic and nonzero for $|q|<1$. \wrt{they are the
products $P_{j}$ of Section~\ref{vmb:sec:accumulation}, there used for
$0<q<1$; they are not the pole sets $\calP_R$ of
Theorem~\ref{vmb:thm:sectors}.}

\begin{lemma}[An analytic two-variable tail]\label{vmb:cw:lem:psi}
For $|q|<1$ define
\begin{equation}\label{vmb:cw:eq:psi}
 \Psi(q,z)=\sum_{k\ge0}zq^k(1-zq^k)
                 \prod_{r=k+1}^{\infty}(1-zq^r)^2.
\end{equation}
The series is analytic in $(q,z)$ locally for $|q|<1$, including near $z=0$,
and
\begin{equation}\label{vmb:cw:eq:tailpsi}
 U_j(q)-U_j^{\le h}(q)=P_j(q)\Psi(q,q^{h+1}).
\end{equation}
Its first coefficients are
\begin{equation}\label{vmb:cw:eq:psifirst}
 \Psi(q,z)=\frac z{1-q}-\frac{z^2}{(1-q)^2}+O(z^3).
\end{equation}
All coefficients are explicitly computable rational functions of $q$.
\end{lemma}
\begin{proof}
On compact $|q|\le r<1$, $|z|\le Z$, the infinite products are uniformly
bounded, while the summands are $O(r^k)$. This proves normal convergence.
Multiply the maximum-part sum for $U_{h+1}$ by
$\prod_{r=h+1}^\infty(1-q^r)^2$ and set $z=q^{h+1}$; the resulting
summands are exactly those in \eqref{vmb:cw:eq:psi}. Equation~\eqref{vmb:cw:eq:Utail}
then proves \eqref{vmb:cw:eq:tailpsi}.

For an explicit coefficient procedure, write
\[
 \prod_{r=1}^\infty(1-zq^r)^2=\sum_{\ell\ge0}\gamma_\ell(q)z^\ell,
 \qquad \gamma_{-1}=0.
\]
Summing the geometric series in $q^{k\ell}$ in \eqref{vmb:cw:eq:psi} gives
\begin{equation}\label{vmb:cw:eq:psicoeff}
 [z^\ell]\Psi(q,z)=\frac{\gamma_{\ell-1}(q)-\gamma_{\ell-2}(q)}{1-q^\ell}
 \quad(\ell\ge1),
\end{equation}
with $\gamma_{-2}=0$ if needed. More explicitly,
\begin{equation}\label{vmb:cw:eq:qbinomcoeff}
 \gamma_\ell(q)=(-1)^\ell\sum_{a+b=\ell}
 \frac{q^{a(a+1)/2+b(b+1)/2}}
 {(q;q)_a(q;q)_b},\qquad
 (q;q)_a=\prod_{r=1}^a(1-q^r).
\end{equation}
For completeness, the coefficient formula for one product follows by
writing $Q(z)=\prod_{r\ge1}(1-zq^r)=(1-zq)Q(qz)$: comparison of coefficients
recursively gives $[z^a]Q=(-1)^aq^{a(a+1)/2}/(q;q)_a$. Squaring proves
\eqref{vmb:cw:eq:qbinomcoeff}. Since $\gamma_0=1$ and $\gamma_1=-2q/(1-q)$,
\eqref{vmb:cw:eq:psifirst} follows.
\end{proof}

\wrt{in this proof $Z$ is a bound and $Q(z)$ a local product; neither is
Part~I's $Z$ or $Q_j$, and the summation indices $a,b$ are local.}

Set $K_j^{\le h}=H_jU_{j+1}^{\le h}$. The lemma implies, locally uniformly
for real $q$ in a compact subinterval of $(0,1)$,
\begin{equation}\label{vmb:cw:eq:generaltail}
 K_j(q)-K_j^{\le h}(q)=a_j(q)q^h+O(q^{2h}),\qquad
 a_j(q)=\frac{q^{j+1}P_{j+1}(q)}{(1-q^j)(1-q)}.
\end{equation}
Each differentiation adds at most a factor of order $h$ to the error
scale. In particular
\begin{equation}\label{vmb:cw:eq:K1tail2}
 K(q)-K^{\le h}(q)
 =P_1(q)q^{h+2}
  -\frac{P_1(q)q^{2h+3}}{1-q}+O(q^{3h}).
\end{equation}

\subsection{The height-cutoff pole to all exponential orders}
Write
\begin{equation}\label{vmb:cw:eq:physicalconst}
 \rho=\rho_1,\quad M_\rho=M_1(\rho),\quad V_\rho=V_1(\rho),\quad
 \alpha=\frac{\rho^2P_1(\rho)}{M_\rho},\quad
 L=\log(1/\rho).
\end{equation}
\wrt{so $M_\rho=M$ and $V_\rho=V$ of \eqref{vmb:eq:MV}, and $L=\log\lambda$
is Part~I's $L$ of \eqref{vmb:eq:interpolant}.}
For sufficiently large integer $h$, let $\widehat\rho_h$ be the root of
$K^{\le h}(\widehat\rho_h)=1$ near $\rho$; it is the dominant pole of $D_1^{\le h}$.
The hat distinguishes a maximum-height cutoff from the minimum-part poles
$\rho_j$. Define
\[
 \delta_h=\log(\widehat\rho_h/\rho).
\]

\begin{theorem}[Height-cutoff exponential expansion]\label{vmb:cw:thm:cutoffseries}
There are polynomials $\varpi_{r-1}(h)$ of degree at most $r-1$ such that, for
every fixed $N\ge1$,
\begin{equation}\label{vmb:cw:eq:cutoffall}
 \delta_h=\sum_{r=1}^{N}\varpi_{r-1}(h)\rho^{rh}
          +O_N(h^N\rho^{(N+1)h}).
\end{equation}
The first two terms are
\begin{align}\label{vmb:cw:eq:cutofftwo}
 \delta_h={}&\alpha\rho^h+
 \left[\alpha^2(h+g)-\frac{\alpha\rho}{1-\rho}\right]\rho^{2h}
 +O(h^2\rho^{3h}),\\
 g={}&2+2\sum_{k\ge1}\frac{k\rho^k}{1-\rho^k}
       -\frac{M_\rho^2+V_\rho}{2M_\rho}.\label{vmb:cw:eq:g}
\end{align}
The leading coefficient of every polynomial is
\begin{equation}\label{vmb:cw:eq:treecoeff}
 [h^{r-1}]\varpi_{r-1}(h)=\frac{r^{r-1}}{r!}\alpha^r.
\end{equation}
Thus the top-degree terms are precisely the expansion of
$-W_0(-\alpha h\rho^h)/h$.
\end{theorem}
\begin{proof}
Let $q=\rho e^\delta$ and $\epsilon=\rho^h$. The root equation is
\[
 K(\rho e^\delta)-1
 =\sum_{r\ge1}A_r(\rho e^\delta)\epsilon^r e^{rh\delta},
\]
where the analytic $A_r$ are obtained from \eqref{vmb:cw:eq:psicoeff}
(Appendix~\ref{vmb:cw:app:cutoff}); in
particular $A_1(q)=q^2P_1(q)$ and
$A_2(q)=-q^3P_1(q)/(1-q)$. At $\delta=0$ the derivative of the left
side is $M_\rho>0$. For small $h\epsilon$, the implicit equation has a
unique solution $\delta=O(\epsilon)$, and its derivative remains bounded
away from zero.

Substitute a formal series in $\epsilon$ and determine its coefficients
successively. At order $r$, the unknown coefficient appears only as
$M_\rho \varpi_{r-1}$. Taylor differentiation of $e^{rh\delta}$ shows
inductively that its remaining right side is a polynomial in $h$ of degree
at most $r-1$. The Taylor residual after $N$ terms is
$O_N(h^N\epsilon^{N+1})$: all derivatives of the analytic $A_r$ and of
$K$ are bounded on a fixed neighborhood, and the exponential factors are
bounded because $h\delta\to0$. Division by the nonvanishing implicit
derivative converts this residual into the claimed root error.

For the first two coefficients, use
\[
 K(\rho e^\delta)-1=M_\rho\delta+
       \frac{M_\rho^2+V_\rho}{2}\delta^2+O(\delta^3),
\]
and $\rho A_1'(\rho)/A_1(\rho)=2+
2\sum k\rho^k/(1-\rho^k)$. This gives \eqref{vmb:cw:eq:cutofftwo}.
To obtain degree $r-1$ in $h$, retain only the linear term on the left
and $A_1(\rho)\epsilon e^{h\delta}$ on the right. Any nonlinear analytic
term, derivative of $A_1$, or $A_r$ with $r\ge2$ lowers the possible degree.
The reduced equation $\delta=\alpha\epsilon e^{h\delta}$ has solution
$-W_0(-\alpha h\epsilon)/h$, whose Lagrange expansion yields
\eqref{vmb:cw:eq:treecoeff}.
\end{proof}

The theorem is an all-orders statement in the ordered scales
$\rho^{rh}$ multiplied by polynomials in $h$. It does not assert convergence
of an unrestricted infinite asymptotic series, nor any Borel or Stokes
property. The small argument of $W_0$ is on its analytic branch at zero.

\begin{writenote}[Added 5 October 2026, batch~98]
The manuscript writes $a_r(q)$ for the order-$r$ coefficients in this proof
and $A_r(q)$ for the same functions in its Appendix~B; they are printed as
$A_r$ in both places, since $a_j(q)$ of \eqref{vmb:cw:eq:generaltail} is a
different family ($a_1=A_1$, but $a_2\ne A_2$). The manuscript's ledger calls
this expansion a ``cutoff-pole transseries''; it is a finite-order expansion
in the scales $h^k\rho^{rh}$ with no convergence, Borel or Stokes claim (as
the paragraph above says), and is not called a transseries here, as for
Theorem~\ref{vmb:thm:W} in Part~I.
\end{writenote}

\section{Tallest columns: a lattice-Gumbel law}
\label{vmb:cw:sec:gumbel}
Let $\mathsf H_n$ be the maximum part of a uniformly chosen unweighted
valley-monotone composition of area $n$.

\begin{theorem}[Uniform maximum-height law]\label{vmb:cw:thm:gumbel}
Uniformly for integers $h=h(n)$ satisfying
$0<\tau_-\le n\rho^h\le\tau_+<\infty$, with fixed $\tau_\pm$,
\begin{equation}\label{vmb:cw:eq:gumbel}
 \boxed{\displaystyle
 \Prob(\mathsf H_n\le h)
 =\exp(-\alpha n\rho^h)
   \left(1+O_{\tau_\pm}\left(\frac{\log n}{n}\right)\right),}
\end{equation}
where $\alpha$ is defined in \eqref{vmb:cw:eq:physicalconst}. Numerically,
\begin{equation}\label{vmb:cw:eq:heightnumbers}
 \begin{gathered}
 \rho=0.5059651577008067\ldots,\qquad
 \alpha=0.4502021724631094\ldots,\\
 L^{-1}=1.4678091752740682\ldots.
 \end{gathered}
\end{equation}
\end{theorem}

\begin{proof}
Choose a fixed radius $R$ strictly between $\rho$ and $\rho_2$, close
enough to $\rho$ that $1-K$ has only its simple root there. Uniform
convergence of $U_2^{\le h}$ to $U_2$ on this disk and Rouch\'e's theorem
show that $1-K^{\le h}$ has exactly one nearby root, the cutoff root
of Theorem~\ref{vmb:cw:thm:cutoffseries}. The higher-level truncated functions
$D_2^{\le h}$ are uniformly analytic on this disk, by the same normally
convergent decomposition used in Lemma~\ref{vmb:cw:lem:poles}. Consequently
\begin{equation}\label{vmb:cw:eq:cappedcoef}
 \coeff qn D_1^{\le h}=\hat c_h\widehat\rho_h^{-n}+O(R^{-n})
\end{equation}
with a uniform error for all sufficiently large $h$.

We record a quantitative estimate for the numerator, rather than assume
it. At a fixed real $q<\rho_2$, subtraction of the two level recurrences
gives an error transported by the multipliers $E_j/(1-K_j)$, whose products
are uniformly bounded for $j\ge2$. The additive kernel perturbation at
level $j$ is bounded by a constant times
\[
 D_{j+1}(q)\,H_j(q)\,
 P_{j+1}(q)\,q^{h+1}=O(q^{2j+h}).
\]
This is summable in $j$, while the final omitted tail has size $O(q^h)$.
It follows that $D_2-D_2^{\le h}=O(q^h)$; differentiating the recurrences
gives $O(hq^h)$ for the first derivative. The same estimates hold locally
with absolute values. Together with the kernel derivative estimates, they
imply
\begin{equation}\label{vmb:cw:eq:residueconv}
 \hat c_h/c_1=1+O(h\rho^h).
\end{equation}
Evaluation at the displaced root does not change this order because
$\widehat\rho_h-\rho=O(\rho^h)$ and $h\rho^h\to0$.

Divide \eqref{vmb:cw:eq:cappedcoef} by \eqref{vmb:cw:eq:basepole} at $j=1$. Theorem
\ref{vmb:cw:thm:cutoffseries} gives
$n\log(\widehat\rho_h/\rho)=\alpha n\rho^h+O(nh\rho^{2h})$.
In the specified window, $h=O(\log n)$ and $\rho^h\asymp n^{-1}$.
Both the exponential error and the residue-ratio error are therefore
$O(\log n/n)$; the uniform Cauchy error is exponentially smaller.
\end{proof}

\begin{writenote}[Added 5 October 2026, batch~98]
The residue-ratio estimate \eqref{vmb:cw:eq:residueconv} is argued in outline
(``differentiating the recurrences gives $O(hq^h)$''); the write records it
as complete in outline, not as an open question. The theorem answers the
second half of Part~I's question ``Joint limits and extreme heights''
(Section~\ref{vmb:sec:future}).
\end{writenote}

\subsection{The lattice phase must not be discarded}
Define
\begin{equation}\label{vmb:cw:eq:center}
 \beta_n=\frac{\log(\alpha n)}L.
\end{equation}
If a subsequence satisfies $\{\beta_n\}\to\theta\in[0,1]$, then, at every
integer $k$,
\begin{equation}\label{vmb:cw:eq:latticephase}
 \Prob(\mathsf H_n-\lfloor \beta_n\rfloor\le k)
 \longrightarrow\exp\left(-\rho^{k-\theta}\right).
\end{equation}
These formulas define lattice distributions; when $\theta=1$, the result
is understood as the one-sided endpoint phase. Different phases generally
give different laws. The centered maximum therefore does not have a single
phase-independent continuous Gumbel limit. Tightness follows by taking
fixed large positive and negative $k$ in \eqref{vmb:cw:eq:latticephase}; in
particular $\mathsf H_n=(\log n)/L+O_{\Prob}(1)$.

\begin{theorem}[Gaussian counts and extremes are independent]\label{vmb:cw:thm:independent}
Put $\mu=1/M_\rho$ and $\sigma^2=V_\rho/M_\rho^3$. Along any phase
subsequence as above,
\begin{equation}\label{vmb:cw:eq:independence}
 \left(\frac{V_n-\mu n}{\sigma\sqrt n},
       \mathsf H_n-\lfloor \beta_n\rfloor\right)
 \Longrightarrow (\chi,Y_\theta),
\end{equation}
where $\chi$ is standard normal, $Y_\theta$ has \eqref{vmb:cw:eq:latticephase}, and
the two limits are independent.
\end{theorem}
\begin{proof}
Mark the valley count by $u=e^{s/\sqrt n}$ and impose the height cutoff.
The untruncated root $r(u)$ is analytic near $u=1$ and differs from $\rho$
by $O(n^{-1/2})$. The truncated-root displacement is, uniformly in this
neighborhood,
\[
 \log\frac{r_h(u)}{r(u)}=
 \frac{u\,a_1(r(u))}{M_1(r(u))}\,r(u)^h
 +O(hr(u)^{2h}).
\]
For $h=\beta_n+O(1)$,
$r(u)^h/\rho^h=1+O(\log n/\sqrt n)$ and the prefactor converges to
$\alpha$. The marked coefficient ratio is therefore the ordinary
quasi-power transform for $V_n$, multiplied by
$\exp(-\alpha n\rho^h)(1+o(1))$. Set $s=it/\sigma$ and include the
centering factor. The first term tends to $e^{-t^2/2}$ and the second tends
to the lattice distribution function. Tightness of both coordinates and
these mixed transforms identify the product limit.
\end{proof}

\wrt{$\mu$, $\sigma^2$ are those of Theorem~\ref{vmb:thm:CLT}, and $r(u)$ is
the root used in its proof.}

\section{Critical extreme heights are Dirichlet mixtures}
\label{vmb:cw:sec:mixed}
The ordinary Gumbel law does not remain unchanged at a critical point.
Several macroscopic regions coexist, and each has its own rate of producing
a column above the cutoff.

Fix $m$ as in Section~\ref{vmb:cw:sec:multi}. Define
\begin{equation}\label{vmb:cw:eq:alphas}
 \alpha_j^*=\frac{u_j^*a_j(\rho_{m+1})}{M_j^*}\quad(1\le j\le m),\qquad
 \alpha_0^*=\frac{a_{m+1}(\rho_{m+1})}{M_{m+1}(\rho_{m+1})},
\end{equation}
where $a_j$ is the explicit product in \eqref{vmb:cw:eq:generaltail}. Index zero
refers to the final $\calD_{m+1}$ block, whose unmarked dominant kernel is
$K_{m+1}$.

\begin{theorem}[Joint critical window and height cutoff]\label{vmb:cw:thm:mixed}
Let $\bm s$ range over a fixed compact subset of $\mathbb R^m$, set
$u_{j,n}=u_j^*e^{s_j/n}$, and let integers $h=h(n)$ satisfy
$0<\tau_-\le\tau_n:=n\rho_{m+1}^h\le \tau_+<\infty$. Then
\begin{align}\label{vmb:cw:eq:mixture}
 &\Prob_{n,\bm u_n}(\mathsf H_n\le h)\notag\\
 &\quad=
 \frac{\displaystyle\int_{\Delta_m}
  \exp\left(\sum_{j=1}^m\frac{s_j\xi_j}{M_j^*}
       -\tau_n\sum_{i=0}^m\alpha_i^*\xi_i\right)\dd\bm \xi}
 {\displaystyle\int_{\Delta_m}
  \exp\left(\sum_{j=1}^m\frac{s_j\xi_j}{M_j^*}\right)\dd\bm \xi}
 +O\left(\frac{\log n}{n}\right).
\end{align}
At exact criticality this is the Laplace transform of
$\sum_{i=0}^m\alpha_i^*X_i$, where
$(X_0,\ldots,X_m)\sim\Dir(1,\ldots,1)$, evaluated at $\tau_n$.
\end{theorem}

\begin{proof}
Truncate every unimodal and valley-monotone factor at height $h$ in
\eqref{vmb:cw:eq:multidecomp}. For a fixed marked level $j\le m$, combine
\eqref{vmb:cw:eq:generaltail} with the implicit derivative at $\rho_{m+1}$. Its nearby
root satisfies
\begin{equation}\label{vmb:cw:eq:mixedroots}
 \log\frac{r_{j,n,h}}{\rho_{m+1}}
 =\alpha_j^*\rho_{m+1}^h-\frac{s_j}{M_j^*n}
 +O\left(h\rho_{m+1}^{2h}+\frac{h\rho_{m+1}^h}{n}+n^{-2}\right).
\end{equation}
Indeed, the derivative of $\log(u_j^*K_j(q))$ with respect to $\log q$
is $M_j^*$; the kernel loss is $u_j^*a_j(\rho_{m+1})\rho_{m+1}^h$ to first order, and the
logarithmic weight increment is $s_j/n$. Taylor's theorem and the
derivative bound following \eqref{vmb:cw:eq:generaltail} give the error.

The truncated terminal function $D_{m+1}^{\le h}$ has a nearby simple pole
whose logarithmic displacement is
$\alpha_0^*\rho_{m+1}^h+O(h\rho_{m+1}^{2h})$. Higher-level denominators are uniformly analytic
on a fixed disk slightly larger than $\rho_{m+1}$. After factoring out these $m+1$
nearby simple poles, the analytic numerator is uniformly bounded on that
disk. At $\rho_{m+1}$ it differs from its critical untruncated value by
$O(h\rho_{m+1}^h+n^{-1})$, by the recurrence estimate in the proof of
Theorem~\ref{vmb:cw:thm:gumbel}. This estimate also handles the possible collision
of some of the moved poles: the analytic factorization does not divide by
their mutual differences.

In the stated window, $\rho_{m+1}^h\asymp n^{-1}$ and $h=O(\log n)$, so
\eqref{vmb:cw:eq:mixedroots} has the form required by Lemma~\ref{vmb:cw:lem:simplex},
with
\[
 \eta_j=\tau_n\alpha_j^*-s_j/M_j^*,\qquad
 \eta_0=\tau_n\alpha_0^*.
\]
Earlier branches in \eqref{vmb:cw:eq:multidecomp} have at most $m-1$ nearby pole
factors and contribute $O(n^{-2})$ relatively (or an exponentially small
amount for $m=1$). Apply the lemma to the dominant branch and divide by
the corresponding untruncated coefficient with $\eta_0=0$ and
$\eta_j=-s_j/M_j^*$. The denominator integral is bounded away from zero
on every compact real $\bm s$-set. This proves \eqref{vmb:cw:eq:mixture}.
\end{proof}

\begin{writenote}[Added 5 October 2026, batch~98]
The manuscript abbreviates $R=R_m$ in this section; every such $R$ is
printed as $\rho_{m+1}$, and its window bounds $a,b$ as $\tau_\pm$. The
intensities $\alpha_j^*$ are not the Lambert coefficients $\alpha_k(c)$ of
\eqref{vmb:eq:alphas}.
\end{writenote}

\subsection{The first critical point in closed form}
When $m=1$, let $\xi_1$ be the macroscopic fraction of area in the bottom-level
renewal blocks. Then $\xi_1$ is uniform on $[0,1]$ at criticality. Writing
$\alpha_1^*$ and $\alpha_0^*$ as in \eqref{vmb:cw:eq:alphas}, the critical-window
height law becomes
\begin{equation}\label{vmb:cw:eq:mixtureone}
 \Prob_{n,u_c e^{s/n}}(\mathsf H_n\le h)
 =\frac{e^{-\tau_n\alpha_0^*}
       G\bigl(s/M_c-\tau_n(\alpha_1^*-\alpha_0^*)\bigr)}{G(s/M_c)}
  +O(\log n/n).
\end{equation}
At $s=0$, this simplifies to
\begin{equation}\label{vmb:cw:eq:twomix}
 \frac{e^{-\tau_n\alpha_1^*}-e^{-\tau_n\alpha_0^*}}
      {\tau_n(\alpha_0^*-\alpha_1^*)}+O(\log n/n),
\end{equation}
with the continuous extension when the two intensities coincide.
The numerical constants are
\begin{align*}
 \alpha_1^*&=0.3053024072567406\ldots,\\
 \alpha_0^*&=0.3765567576528582\ldots,\\
 \frac1{\log(1/\rho_2)}&=2.1533787748784431\ldots.
\end{align*}
Thus critical maximum height has scale $2.1533\ldots\log n$, compared with
$1.4678\ldots\log n$ in the ordinary unweighted model.

The limiting expression is not a single exponential $e^{-c\tau}$ unless
all active intensities coincide. Indeed, the second Taylor coefficient
would otherwise force
$\E(\sum\alpha_i^*X_i)^2=(\E\sum\alpha_i^*X_i)^2$, which fails for a
nonconstant linear function on the simplex. This distinguishes an actual
mixture from merely replacing the Gumbel intensity by its mean.

For the lattice statement, take
$\beta_n^{\mathrm{crit}}=\log n/\log(1/\rho_{m+1})$ and a subsequence with fractional
parts tending to $\theta$. At integer $k$, put
$\tau=\rho_{m+1}^{k-\theta}$ in \eqref{vmb:cw:eq:mixture}. The resulting distribution
function is the phase-dependent Dirichlet mixture. Joint marking of the
area fractions shows more: conditional on their limiting values $X_i$,
the extreme-height distribution has intensity $\sum\alpha_i^*X_i$.
Consequently the maximum is generally \emph{not} asymptotically independent
of the critical macroscopic fractions, unlike the Gaussian independence
in Theorem~\ref{vmb:cw:thm:independent}.

\begin{proposition}[{[write]} Conditional intensity of the critical maximum]
\label{vmb:cw:prop:conditional}
In the setting of Theorem~\ref{vmb:cw:thm:mixed}, let $\bm X_n$ be the vector
in \eqref{vmb:cw:eq:dirichlet} on the dominant branch (and $\bm X_n=0$ on the
earlier branches), and let $\pi_{\bm s}$ be the probability measure on the
simplex with density proportional to $\exp(\sum_{j=1}^ms_j\xi_j/M_j^*)$.
Along any subsequence with $\tau_n\to\tau\in[\tau_-,\tau_+]$, for every
continuous function $f$ on $[0,1]^{m+1}$,
\begin{equation}\label{vmb:cw:eq:conditional}
 \E_{n,\bm u_n}\bigl[f(\bm X_n);\,\mathsf H_n\le h\bigr]
 \longrightarrow\int f(\bm\xi)\,
 \exp\Bigl(-\tau\sum_{i=0}^m\alpha_i^*\xi_i\Bigr)\,\pi_{\bm s}(\dd\bm\xi).
\end{equation}
Thus, in the limit, $\bm X_n$ has law $\pi_{\bm s}$ and, conditionally on
$\bm X=\bm\xi$, the event $\{\mathsf H_n\le h\}$ has probability
$\exp(-\tau\sum_i\alpha_i^*\xi_i)$. The limits of $\bm X_n$ and of
$\mathbf 1\{\mathsf H_n\le h\}$ are independent if and only if
$\alpha_0^*=\alpha_1^*=\cdots=\alpha_m^*$.
\end{proposition}
\begin{proof}
This proof was added in writing; the manuscript asserts the statement in
the paragraph above. Mark areas as in the proof of
Theorem~\ref{vmb:cw:thm:dirichlet}, by $v_j/n$ in the level-$j$ repeated units
and $v_0/n$ in the final block, in the truncated decomposition of the proof
of Theorem~\ref{vmb:cw:thm:mixed}. Replacing $q$ by $qe^{v_j/n}$ in
$u_{j,n}K_j^{\le h}$ multiplies its root $r_{j,n,h}$ by exactly
$e^{-v_j/n}$, and likewise for the cutoff pole of $D_{m+1}^{\le h}$; the
remaining analytic factors, evaluated at $qe^{v/n}$, are uniformly bounded
and change by $O(1/n)$ at $\rho_{m+1}$. Lemma~\ref{vmb:cw:lem:simplex} applies with
$\eta_j=\tau_n\alpha_j^*-s_j/M_j^*-v_j$ and $\eta_0=\tau_n\alpha_0^*-v_0$, and
the earlier branches contribute $O(n^{-2})$ relatively, as in that proof.
Dividing by the unmarked untruncated coefficient gives, uniformly for
$\bm v$ in compact subsets of $\mathbb R^{m+1}$,
\[
 \E_{n,\bm u_n}\bigl[e^{\langle\bm v,\bm X_n\rangle};\,\mathsf H_n\le h\bigr]
 =\int e^{\langle\bm v,\bm\xi\rangle}
 e^{-\tau_n\langle\bm\alpha^*,\bm\xi\rangle}\,\pi_{\bm s}(\dd\bm\xi)
 +O\Bigl(\frac{\log n}n\Bigr),
\]
with $\langle\bm v,\bm\xi\rangle=\sum_{i=0}^mv_i\xi_i$. The left side is the
Laplace transform of the finite measure
$\mu_n(B)=\Prob_{n,\bm u_n}(\bm X_n\in B,\ \mathsf H_n\le h)$ on the compact
set $[0,1]^{m+1}$. As $\tau_n\to\tau$, these transforms converge for every
real $\bm v$ to that of
$\mu(\dd\bm\xi)=e^{-\tau\langle\bm\alpha^*,\bm\xi\rangle}\pi_{\bm s}(\dd\bm\xi)$,
whose total mass is positive. Hence the total masses converge, and the
normalized probability measures, whose moment generating functions converge
everywhere, converge weakly; so $\mu_n\to\mu$ weakly, which is
\eqref{vmb:cw:eq:conditional}. With Proposition~\ref{vmb:cw:prop:tilted}
($\bm X_n\Rightarrow\pi_{\bm s}$) this identifies the joint limit as stated.
It is a product law exactly when $e^{-\tau\langle\bm\alpha^*,\bm\xi\rangle}$
is $\pi_{\bm s}$-almost surely constant. Since $\pi_{\bm s}$ has a positive
continuous density on the simplex and
$\langle\bm\alpha^*,\bm\xi\rangle=\alpha_0^*+\sum_{j=1}^m(\alpha_j^*-\alpha_0^*)\xi_j$
there, with $\tau>0$, this happens if and only if all $\alpha_i^*$ are equal.
\end{proof}

\wrt{at the first critical point the numerical values
$\alpha_1^*=0.3053\ldots$ and $\alpha_0^*=0.3766\ldots$ above differ, so the
critical maximum is not asymptotically independent of the area fraction
there; these decimals are diagnostics, not interval certificates.}

\section{Inversion of the new asymptotic expansions}
\label{vmb:cw:sec:inversion}

\subsection{Inverting the height-cutoff pole}\label{vmb:cw:sec:inverseheight}
A discrete cutoff has no automatic real inverse. Here an exact analytic
interpolation is available: for large real $h$ and real $q$ near $\rho$,
define
\begin{equation}\label{vmb:cw:eq:heightinterp}
 \widehat K_h(q)=K(q)-H_1(q)P_2(q)\Psi(q,q^{h+1}).
\end{equation}
It agrees with the height-truncated kernel at every sufficiently large
integer $h$. Its nearby root defines $\delta(h)$ as in
Theorem~\ref{vmb:cw:thm:cutoffseries}. The derivative of the leading tail with
respect to $h$ has the sign of $\log\rho<0$, while subsequent derivatives
are smaller, so $\delta(h)$ is strictly decreasing for large $h$.

\begin{lemma}[{[write]} The interpolated shift decreases]\label{vmb:cw:lem:deltamono}
For all sufficiently large real $h$, $\widehat K_h(q)=1$ has a unique root
$q=\rho e^{\delta(h)}$ with $\delta(h)$ small; $\delta$ is real analytic,
$\delta(h)=\alpha\rho^h(1+O(h\rho^h))$, and $\delta'(h)<0$. At integers
$h\ge2$, $\widehat K_h=K^{\le h}$ and $\delta(h)=\delta_h$.
\end{lemma}
\begin{proof}
This proof was added in writing, expanding the one-line justification
above. By \eqref{vmb:cw:eq:tailpsi} with $j=2$, $K-K^{\le h}=H_1(U_2-U_2^{\le h})
=H_1P_2\Psi(q,q^{h+1})$ for integers $h\ge2$, which gives the agreement.
Since $\Psi(q,z)=\sum_{\ell\ge1}\psi_\ell(q)z^\ell$ converges for small $|z|$,
locally uniformly in $q$, we have
$K(q)-\widehat K_h(q)=\sum_{r\ge1}A_r(q)q^{rh}$ with the $A_r$ of
Appendix~\ref{vmb:cw:app:cutoff}, for all real $h$ large. Put
$\Phi(\delta,h)=\widehat K_h(\rho e^\delta)-1
=K(\rho e^\delta)-1-\sum_{r\ge1}A_r(\rho e^\delta)\rho^{rh}e^{rh\delta}$,
analytic in $(\delta,h)$ for $|\delta|$ small and $h$ large. On
$|\delta|\le1/h^2$, say, $rh|\delta|\le r/h$, so
$\partial_\delta\Phi=M_\rho+O(h\rho^h)>0$ and the implicit function theorem
gives a unique small root $\delta(h)$, analytic in $h$, with
$\delta(h)=\alpha\rho^h(1+O(h\rho^h))$ by the first-order balance
$M_\rho\delta=A_1(\rho)\rho^h+\cdots$ and $A_1(\rho)/M_\rho=\alpha$. Next,
$\partial_h\Phi=-\sum_{r\ge1}r(\log\rho+\delta)A_r(\rho e^\delta)
\rho^{rh}e^{rh\delta}=L\,A_1(\rho)\rho^h(1+O(h\rho^h))>0$, because
$A_1(\rho)=\rho^2P_1(\rho)>0$ and $L=-\log\rho>0$. Hence
$\delta'(h)=-\partial_h\Phi/\partial_\delta\Phi<0$.
\end{proof}

\begin{theorem}[Inverse exponential expansion]\label{vmb:cw:thm:inverseheight}
Let $h(\delta)$ be the large real inverse of \eqref{vmb:cw:eq:heightinterp},
where $\delta\downarrow0$, and put
\begin{equation}\label{vmb:cw:eq:h0}
 h_0=\frac{\log(\alpha/\delta)}L.
\end{equation}
Then
\begin{equation}\label{vmb:cw:eq:inverseheight}
 \boxed{\displaystyle
 h(\delta)=h_0+\frac\delta L
 \left(h_0+g-\frac{\rho}{\alpha(1-\rho)}\right)
 +O(\delta^2h_0^2).}
\end{equation}
To every order it has an expansion in powers of $\delta$ whose coefficients
are polynomials in $h_0$; the coefficient at $\delta^r$ has degree at most
$r$.
\end{theorem}
\begin{proof}
Set $\epsilon=\rho^h$. Formula~\eqref{vmb:cw:eq:cutofftwo} is
\[
 \delta=\alpha\epsilon+
 \left[\alpha^2(h+g)-\frac{\alpha\rho}{1-\rho}\right]\epsilon^2
 +O(h^2\epsilon^3).
\]
Solving first for $\epsilon$ and taking its logarithm gives
\[
 -Lh=\log(\delta/\alpha)
 -\delta\left(h+g-\frac{\rho}{\alpha(1-\rho)}\right)
 +O(\delta^2h^2).
\]
Here $h=h_0+O(\delta h_0)$, so substitution of $h_0$ in the first correction
proves \eqref{vmb:cw:eq:inverseheight}. Repeating the same substitution with
\eqref{vmb:cw:eq:cutoffall} gives the asserted polynomial structure at each finite
order. The derivative of the logarithmic equation with respect to $h$ is
$-L+O(\delta h)$, bounded away from zero, so every formal residual yields
a corresponding asymptotic remainder. At order $N$, the remainder is
$O_N(\delta^{N+1}h_0^{N+1})$.
\end{proof}

At integer cutoffs, a truncated inverse must be accompanied by a rounding
interval. The exact interpolation supplies a well-defined crossing; a
finite expansion supplies only an approximation to it. Its ceiling is
unambiguous only when the approximation is separated from the integers by
more than the proved remainder.

\subsection{Critical counting inversion is almost exactly Lambert--\texorpdfstring{$W$}{W}}
Let $L_c=\log(1/\rho_2)$ and $a_c=\kappa_0/\kappa$. The entire dominant-pole
polynomial in \eqref{vmb:cw:eq:criticalcount} has the real interpolation
\[
 Y_c(x)=\kappa(x+a_c)e^{L_cx}.
\]
For sufficiently large $Y$ its inverse is exactly
\begin{equation}\label{vmb:cw:eq:Wcrit}
 \boxed{\displaystyle
 x_c(Y)=\frac1{L_c}W_0\left(\frac{L_c}{\kappa}Ye^{L_ca_c}\right)-a_c.}
\end{equation}
This follows by putting $z=L_c(x+a_c)$ in $Y=Y_c(x)$. No additional
inverse-power corrections to this specified interpolant are missing:
all errors in replacing the actual coefficients by $Y_c(n)$ are
exponentially smaller.

For example, with $z_c=(L_c/\kappa)Ye^{L_ca_c}$, $\ell_1=\log z_c$, and
$\ell_2=\log \ell_1$, ordinary series substitution in $W+\log W=\log z_c$ gives
\begin{equation}\label{vmb:cw:eq:Wlog}
 x_c(Y)=\frac1{L_c}\left(
 \ell_1-\ell_2+\frac{\ell_2}{\ell_1}
 +\frac{\ell_2(\ell_2-2)}{2\ell_1^2}
 +O\left(\frac{\ell_2^3}{\ell_1^3}\right)\right)-a_c.
\end{equation}
The formula \eqref{vmb:cw:eq:Wcrit} is preferable numerically; \eqref{vmb:cw:eq:Wlog}
exhibits the inverse-logarithmic scale.

\begin{proposition}[Discrete crossing bounds]\label{vmb:cw:prop:rounding}
For $Y>0$, define $N_c(Y)=\min\{n\ge1:f_n(u_c)\ge Y\}$. For some
constant $C_*>0$, with $R$ as in \eqref{vmb:cw:eq:criticalcount},
\begin{equation}\label{vmb:cw:eq:rounding}
 \left\lceil x_c(Y)-C_*\frac{R^{-x_c(Y)}}Y\right\rceil
 \le N_c(Y)\le
 \left\lceil x_c(Y)+C_*\frac{R^{-x_c(Y)}}Y\right\rceil
\end{equation}
for all sufficiently large $Y$.
\end{proposition}
\begin{proof}
Appending a part one creates no new valley and preserves all existing
valley weights. It gives a weight-preserving injection from area $n$ to
area $n+1$; the singleton $(n+1)$ is an additional object. Thus $f_n(u_c)$
is strictly increasing. Near $x_c(Y)$, $Y_c'(x)\asymp Y$. The coefficient
error $O(R^{-n})$ therefore moves the real crossing by at most a constant
times $R^{-x_c(Y)}/Y$. A mean-value bound on a fixed neighborhood of the
crossing proves the two inequalities; outside that neighborhood the leading
exponential already separates the values. The displayed uncertainty tends
to zero because $R>\rho_2$.
\end{proof}

\subsection{Multicritical counting inverses}
At an $m$-level critical point, write
$\Pi_m(x)=\kappa^{(m)}x^m(1+\pi_1/x+\cdots+\pi_m/x^m)$ and
$L_m=\log(1/\rho_{m+1})$ \wrt{so $L_1=L_c$}. The leading inverse is
\begin{equation}\label{vmb:cw:eq:Wmulti}
 x_0(Y)=\frac m{L_m}W_0\left(
       \frac{L_m}{m}(Y/\kappa^{(m)})^{1/m}\right).
\end{equation}
The inverse of $e^{L_mx}\Pi_m(x)$ is obtained to all algebraic orders by
setting $x=x_0+\varphi$ in
\begin{equation}\label{vmb:cw:eq:multiinverseimplicit}
 \varphi=-\frac1{L_m}\left[
 m\log(1+\varphi/x_0)+
 \log\left(1+\frac{\pi_1}{x_0+\varphi}+\cdots+
                    \frac{\pi_m}{(x_0+\varphi)^m}\right)\right].
\end{equation}
The right side and its derivative are $O(x_0^{-1})$, so local analytic
Lagrange inversion is convergent for large $x_0$ and gives an effective
coefficient algorithm. In particular,
\[
 x=x_0-\frac{\pi_1}{L_mx_0+m}+O(x_0^{-2}).
\]
The exponentially small coefficient remainder and the same injection give
rounding bounds of the form \eqref{vmb:cw:eq:rounding}, with the corresponding
polynomial interpolant and radius. These inversion tools are standard;
the new input is the explicitly identified critical polynomial and its
phase-dependent degree.

\subsection{Height quantiles and their rounding phase}
For fixed $p\in(0,1)$, define the integer quantile
$h_p(n)=\min\{h:\Prob(\mathsf H_n\le h)\ge p\}$. The real leading crossing is
\begin{equation}\label{vmb:cw:eq:heightquantile}
 \beta_p(n)=\frac{\log(\alpha n)-\log(-\log p)}L.
\end{equation}
Theorem~\ref{vmb:cw:thm:gumbel} and the nonzero derivative of
$\exp(-\alpha n\rho^x)$ at $x=\beta_p(n)$ imply, for some $C_p$,
\begin{equation}\label{vmb:cw:eq:quantilerounding}
 \left\lceil \beta_p(n)-C_p\frac{\log n}{n}\right\rceil
 \le h_p(n)\le
 \left\lceil \beta_p(n)+C_p\frac{\log n}{n}\right\rceil.
\end{equation}
A bare ceiling of \eqref{vmb:cw:eq:heightquantile} would omit the narrow but real
rounding ambiguity. At criticality, replace the exponential by the monotone
mixture Laplace transform in \eqref{vmb:cw:eq:mixture}; its unique positive real
crossing in $\tau$ supplies the analogous quantile constant.

\begin{proposition}[{[write]} Multicritical rounding and critical quantiles]
\label{vmb:cw:prop:roundingmulti}
\textup{(i)} At the weights \eqref{vmb:cw:eq:multiweights}, let
$f^{(m)}_n=\coeff qn F_m(q,\bm u^*)$, $N_m(Y)=\min\{n\ge1:f^{(m)}_n\ge Y\}$,
and let $x_m(Y)$ be the inverse of $Y_m(x)=e^{L_mx}\Pi_m(x)$ on its
eventually increasing branch. With $R$ as in \eqref{vmb:cw:eq:multicount},
there is $C_*>0$ such that, for all sufficiently large $Y$,
\[
 \left\lceil x_m(Y)-C_*\frac{R^{-x_m(Y)}}Y\right\rceil
 \le N_m(Y)\le
 \left\lceil x_m(Y)+C_*\frac{R^{-x_m(Y)}}Y\right\rceil.
\]
\textup{(ii)} In the setting of Theorem~\ref{vmb:cw:thm:mixed} with fixed
$\bm s$, let $\Phi_{\bm s}(\tau)$ be the main term of
\eqref{vmb:cw:eq:mixture} as a function of $\tau=\tau_n$. For fixed
$p\in(0,1)$ it has a unique root $\tau_p>0$ of $\Phi_{\bm s}(\tau)=p$; put
$\beta^{\mathrm{crit}}_p(n)=(\log n-\log\tau_p)/L_m$. Then
$h_p(n)=\min\{h:\Prob_{n,\bm u_n}(\mathsf H_n\le h)\ge p\}$ satisfies, for
some $C_p$ and all large $n$,
\[
 \left\lceil \beta^{\mathrm{crit}}_p(n)-C_p\frac{\log n}{n}\right\rceil
 \le h_p(n)\le
 \left\lceil \beta^{\mathrm{crit}}_p(n)+C_p\frac{\log n}{n}\right\rceil.
\]
The same argument, with $\Phi(\tau)=e^{-\alpha\tau}$ and
Theorem~\ref{vmb:cw:thm:gumbel}, proves \eqref{vmb:cw:eq:quantilerounding}.
\end{proposition}
\begin{proof}
This proof was added in writing; the manuscript states (i) and (ii) by
analogy in the two paragraphs above.
(i) Appending a part one to a valley-monotone composition creates no
valley (the new last run is an endpoint run, and the run before it gains a
smaller right neighbour, so it is not a valley) and removes none, so it is a
weight-preserving injection from area $n$ to area $n+1$; the singleton
$(n+1)$, of weight one, is not in its image. Hence $f^{(m)}_{n+1}>f^{(m)}_n$.
By \eqref{vmb:cw:eq:multicount}, $f^{(m)}_n=Y_m(n)+O(R^{-n})$, and
$Y_m'(x)=Y_m(x)(L_m+\Pi_m'(x)/\Pi_m(x))\asymp Y_m(x)$ for large $x$. The proof
of Proposition~\ref{vmb:cw:prop:rounding} now applies verbatim with $Y_m$, $x_m$
in place of $Y_c$, $x_c$.

(ii) With $S=\sum_i\alpha_i^*\xi_i$, $\Phi_{\bm s}(\tau)$ is the expectation
of $e^{-\tau S}$ under the tilted law $\pi_{\bm s}$ of
Proposition~\ref{vmb:cw:prop:conditional}. Since $S\ge\min_i\alpha_i^*>0$,
$\Phi_{\bm s}$ is continuous, $\Phi_{\bm s}(0)=1$, $\Phi_{\bm s}(\tau)\to0$
as $\tau\to\infty$, and $\Phi_{\bm s}'(\tau)=-\E_{\pi_{\bm s}}[Se^{-\tau S}]<0$;
so $\tau_p$ exists and is unique. Put $g(x)=\Phi_{\bm s}(n\rho_{m+1}^x)$ for
real $x$, so $g(\beta^{\mathrm{crit}}_p(n))=p$ and
$g'(x)=-L_m\tau\Phi_{\bm s}'(\tau)$ at $\tau=n\rho_{m+1}^x$, which is
$\ge c>0$ while $\tau\in[\rho_{m+1}^2\tau_p,\tau_p/\rho_{m+1}^2]$. Apply
Theorem~\ref{vmb:cw:thm:mixed} with $\tau_-=\rho_{m+1}^2\tau_p$,
$\tau_+=\tau_p/\rho_{m+1}^2$; its error is at most $K\log n/n$ there. Let
$\varepsilon_n=C_p\log n/n$ with $C_p>K/c$. The integer
$h^+=\lceil\beta^{\mathrm{crit}}_p+\varepsilon_n\rceil$ has $\tau$ in that
range and $g(h^+)\ge p+c\varepsilon_n$, so
$\Prob(\mathsf H_n\le h^+)\ge p+c\varepsilon_n-K\log n/n>p$ and
$h_p(n)\le h^+$. The integer $h^-=\lceil\beta^{\mathrm{crit}}_p-\varepsilon_n\rceil-1$
lies in $[\beta^{\mathrm{crit}}_p-\varepsilon_n-1,\beta^{\mathrm{crit}}_p-\varepsilon_n)$,
so its $\tau$ is in range and $g(h^-)<p-c\varepsilon_n$; hence
$\Prob(\mathsf H_n\le h)\le\Prob(\mathsf H_n\le h^-)<p$ for every $h\le h^-$,
and $h_p(n)\ge h^-+1$. These are the two ceilings.
\end{proof}

\begin{writenote}[Added 5 October 2026, batch~98: instances of repository results]
As for Part~I's inverse (the notes in Sections~\ref{vmb:sec:accumulation}
and~\ref{vmb:sec:inverse}), the inversion tools of this section are
instances of the transseries-and-inversion volume~\cite{vmb-tai}, and the
manuscript itself calls them standard. Writing $X=x+a_c$,
$Y=Y_c(x)$ reads $L_cX+\log X=\log(Y/\kappa)+L_ca_c$, the Lambert core
\texttt{p0:thm:lambert-core} with $a=L_c$, $b=1>0$ (principal branch); this
is \eqref{vmb:cw:eq:Wcrit}, and \eqref{vmb:cw:eq:Wlog} is its standard
logarithmic series. At $m$ levels the leading block
$L_mx+m\log x=\log(Y/\kappa^{(m)})$ is the same core with $a=L_m$, $b=m$, giving
\eqref{vmb:cw:eq:Wmulti}, and \eqref{vmb:cw:eq:multiinverseimplicit} is
reversion about that core (\texttt{p0:thm:core-reversion}), i.e.\ Lagrange
inversion (\texttt{p0:thm:LB}), which also underlies
Theorem~\ref{vmb:cw:thm:inverseheight}. The reduced equation
$\delta=\alpha\epsilon e^{h\delta}$ in the proof of
Theorem~\ref{vmb:cw:thm:cutoffseries} is solved by the Lambert function in the
same way. Propositions~\ref{vmb:cw:prop:rounding}
and~\ref{vmb:cw:prop:roundingmulti}(i), and \eqref{vmb:cw:eq:quantilerounding},
are the staircase theorem \texttt{p0:thm:staircase}, parts (1)--(2), whose
generic arithmetic is formalized in Lean as \texttt{Fabius.staircase\_ceil}
and \texttt{Fabius.staircase\_separation} in
\nolinkurl{Analysis/FabiusFunction/Lean/FabiusFunction/StaircaseInversion.lean};
those statements concern an arbitrary monotone interpolation and give no
formal status to the propositions here. No novelty is claimed for the
methods; the constants and the critical polynomials are specific to this
report.
\end{writenote}

\section{Reproducible checks and numerical behavior}
\label{vmb:cw:sec:checks}

\subsection{Exact finite verification}
The supplied \texttt{code/verify.py} uses only Python's standard library.
It computes coefficients by triangular formal-series division over the
integers. Independently it generates every composition through area fourteen,
extracts plateau valleys directly, and compares six weighted or
height-restricted settings. All $84$ comparisons pass. The settings include
unweighted objects, maximum heights three and five, weight two at level one,
and simultaneous weights two and three at levels one and two. It also checks
A001523 through area twenty and retains unweighted coefficients through area
$240$.

The exact denominator for a marked level is
\begin{equation}\label{vmb:cw:eq:exactden}
 Q_{j,u}(q)=(1-q^j)^2-u q^j(1-q^j)U_{j+1}(q),
\end{equation}
and the numerator in $Q_{j,u}D_j$ is
\begin{equation}\label{vmb:cw:eq:exactnum}
 q^j(1-q^j)-u q^{2j}U_{j+1}(q)+D_{j+1}(q).
\end{equation}
Because $Q_{j,u}(0)=1$, no division of integers is needed. At height cap $h$,
initialize $U_{h+1}=D_{h+1}=0$; without a cap, initializing at level $N+1$
is exact through degree $N$. These are finite identities, not asymptotic
approximations.

\begin{writenote}[Added 5 October 2026, batch~98]
$Q_{j,u}=(1-q^j)(1-q^j-uq^jU_{j+1})$, so $Q_{j,1}=(1-q^j)Q_j$ with Part~I's
$Q_j$ of \eqref{vmb:eq:HKQ}, and \eqref{vmb:cw:eq:exactden}--\eqref{vmb:cw:eq:exactnum}
at $u=1$ are the last line of \eqref{vmb:eq:algorithm}.
\end{writenote}

\subsection{Critical coefficients and the double pole}
The numerical programs use $65$-digit \texttt{mpmath} arithmetic with a
maximum level of $320$. Comparing levels $220$ and $320$ in $K(\rho_2)$ gives
a difference about $1.43\cdot10^{-43}$. This is a convergence diagnostic,
not a rigorous interval certificate. The exact critical constants are
always the analytic quantities in \eqref{vmb:cw:eq:critical}--\eqref{vmb:cw:eq:kappas}.

\begin{center}
\begin{tabular}{@{}rrrr@{}}
\toprule
$n$ & $\rho_2^nf_n(u_c)-(\kappa n+\kappa_0)$ & $\E V_n/n$ & $\Var V_n/n^2$\\
\midrule
30  & $4.6043\cdot10^{-2}$ & $0.02736564$ & $0.00084894$\\
60  & $2.2957\cdot10^{-3}$ & $0.03149433$ & $0.00065425$\\
120 & $5.5586\cdot10^{-6}$ & $0.03413759$ & $0.00055328$\\
240 & $3.2472\cdot10^{-11}$& $0.03564653$ & $0.00050604$\\
\midrule
Limit & $0$ & $0.03728019$ & $0.00046327$\\
\bottomrule
\end{tabular}
\end{center}
The exponentially small residual concerns the coefficient formula. The
moments need not converge exponentially: marking splits the double pole,
and their leading finite-size corrections are algebraic.

At $n=240$ the leading critical-window formula, its first correction, and
the independently computed scaled coefficient are:
\begin{center}
\begin{tabular}{@{}rrrr@{}}
\toprule
$s$ & $\kappa nG(s/M_c)$ & with $G_1(s)$ & $\rho_2^nf_n(u_c e^{s/n})$\\
\midrule
$-4$ & $48.27503236$ & $48.32923935$ & $48.31748024$\\
$0$  & $55.83114319$ & $55.49838473$ & $55.49838473$\\
$4$  & $65.04998792$ & $64.24315456$ & $64.26440169$\\
\bottomrule
\end{tabular}
\end{center}
The remaining error for nonzero $s$ is consistent with the proved $O(1/n)$
scaled remainder; the table is not used to infer that bound.

\subsection{Extreme heights and pole shifts}
For the ordinary model, the exact finite coefficient ratio and the leading
height prediction in \eqref{vmb:cw:eq:gumbel} give:
\begin{center}
\begin{tabular}{@{}rrrr@{}}
\toprule
$n$ & $h$ & exact $\Prob(\mathsf H_n\le h)$ & $e^{-\alpha n\rho^h}$\\
\midrule
60  &5& $0.39227262$ &$0.40832225$\\
120 &6& $0.39622717$ &$0.40398216$\\
240 &7& $0.39582494$ &$0.39963725$\\
\bottomrule
\end{tabular}
\end{center}
The visible finite-size differences are retained rather than hidden by a
fitted intensity. The pole-shift expansion itself converges much faster:
\begin{center}
\begin{tabular}{@{}rrr@{}}
\toprule
$h$ & relative error, first term & relative error, two terms\\
\midrule
8  & $-1.7175\cdot10^{-2}$ & $-4.3993\cdot10^{-4}$\\
12 & $-1.6241\cdot10^{-3}$ & $-3.9517\cdot10^{-6}$\\
16 & $-1.3958\cdot10^{-4}$ & $-2.9209\cdot10^{-8}$\\
24 & $-8.8483\cdot10^{-7}$ & $-1.1741\cdot10^{-12}$\\
32 & $-5.0260\cdot10^{-9}$ & $-3.7886\cdot10^{-17}$\\
\bottomrule
\end{tabular}
\end{center}
Here ``relative error'' means approximation divided by the computed
$\delta_h$, minus one.

At the first critical point, the numerically evaluated finite weighted
coefficient ratios and the
mixture prediction \eqref{vmb:cw:eq:twomix} include
\begin{center}
\begin{tabular}{@{}rrrr@{}}
\toprule
$n$ & $h$ & computed critical CDF & mixture prediction\\
\midrule
120&7 &$0.19320378$&$0.20584639$\\
120&8 &$0.37460867$&$0.36990049$\\
240&9 &$0.28593991$&$0.28662767$\\
240&10&$0.46399389$&$0.45564238$\\
\bottomrule
\end{tabular}
\end{center}
These computations use the explicit weighted recurrence, not samples from
a Monte Carlo generator.

Finally, at the two-level critical point,
\[
 \rho_3=0.6949152510\ldots,\quad
 u_1^*=0.02240337599\ldots,\quad
 u_2^*=0.21458477859\ldots.
\]
For $n=240$ and $(s_1,s_2)=(3,-2)$, the numerically evaluated coefficient ratio with the
$1/n$ weight tilts is $1.0118132157\ldots$, whereas the limiting simplex
integral is $1.0129148706\ldots$. For $(-3,2)$ the corresponding numbers
are $0.9911396048\ldots$ and $0.9897390369\ldots$. These are additional
finite checks of the multilevel construction, not substitutes for
Theorem~\ref{vmb:cw:thm:dirichlet}.

\subsection{Running the package}
From the package directory:
\begin{verbatim}
python code/verify.py
python code/symbolic_checks.py
python code/numerics.py
python code/critical_extremes.py
pdflatex article.tex
pdflatex article.tex
\end{verbatim}
The first script requires only Python's standard library; symbolic checks
require \texttt{sympy}, and numerical diagnostics require \texttt{mpmath}.
The data directory records the outputs. No external papers, repository
source files, or font files are included. The scripts do not change the
remote repository or submit anything to OEIS.

\begin{writenote}[Added 5 October 2026, batch~98]
The commands above describe the delivered package, not this repository. The
four programs are shipped in \texttt{code/} and their outputs in
\texttt{data/}, each under its delivered name with the prefix
\texttt{02-critical-windows-} (for example
\texttt{code/\allowbreak 02-critical-windows-\allowbreak verify.py}); the
delivered article, its PDF and
its build script are not shipped (this report's README gives the build
command). The programs write unprefixed files into a sibling \texttt{data/}
directory, \texttt{numerics.py} and \texttt{critical\_extremes.py} import
\texttt{verify} and \texttt{numerics} by their delivered names (inside the
repository those names would resolve to Part~I's programs, which lack the
imported functions), and \texttt{critical\_extremes.py} reads
\texttt{numerical\_checks.json}. The README therefore reruns them on a scratch
copy with the delivered names restored. The rerun of the write (5~October
2026, Python 3.13.5, SymPy 1.14.0, mpmath 1.3.0) passed: 84 enumeration checks
and A001523 through area~20, five symbolic identity groups, and the two
numerical programs, all four outputs equal to the shipped files as JSON (the
bytes differ only by CRLF line ends written on Windows). In the data keys,
\texttt{R} is $\rho_2$, \texttt{R3} is $\rho_3$, \texttt{P} is $P_1(\rho)$,
\texttt{A} is $\alpha$, \texttt{physical\_M}, \texttt{physical\_V} are
Part~I's $M$, $V$, \texttt{critical\_M}, \texttt{critical\_V},
\texttt{Jc}, \texttt{critical\_D}, \texttt{critical\_d} are $M_c$, $V_c$,
$J_c$, $D_c$, $d_c$, \texttt{critical\_height\_alpha1} and
\texttt{critical\_height\_alpha2} are $\alpha_1^*$ and $\alpha_0^*$,
\texttt{critical\_level2\_M} is $M_2(\rho_2)$, \texttt{shift} is $a_c$, and
\texttt{u1}, \texttt{u2}, \texttt{M1}, \texttt{M2} of the two-level record are
$u_1^*$, $u_2^*$, $M_1^*$, $M_2^*$ at $m=2$. Keys such as
\texttt{exact\_cdf} and \texttt{exact\_ratio} denote the finite recurrence
evaluated at approximate weights, not enclosures. The A001523 terms through
area~20 hard-coded in \texttt{02-critical-windows-verify.py} are OEIS data
(CC~BY-SA~4.0); the arrays of $D_1$, $U_2$, $D_2$ in
\texttt{02-critical-windows-exact\_checks.json} are computed, not OEIS data.
\end{writenote}

\section{Further questions and research}
\label{vmb:cw:sec:future}

\begin{writenote}[Added 5 October 2026, batch~98]
Subsections~\ref{vmb:cw:sec:q-pointprocess}--\ref{vmb:cw:sec:q-formal} are
the manuscript's Section~11 (``Further research questions''), printed with
its numbering of topics 1--7. Subsection~\ref{vmb:cw:sec:q-write} lists, under
the repository's standing rule of 4~October 2026, every statement of the
manuscript that it supports only by a sketch, an analogy or a one-line
remark, with what the write did about each, and the part of Part~I's question
that Part~II leaves open.
\end{writenote}

\subsection{A critical extremal point process}
\label{vmb:cw:sec:q-pointprocess}
Theorem~\ref{vmb:cw:thm:mixed} identifies the law of the largest column. A stronger
next target is the joint distribution of the largest several column heights,
with their plateau multiplicities and the levels of their renewal components.
A natural conjecture is a component-marked Poisson limit conditional on the
Dirichlet fractions, with a corresponding random-intensity mixture after
conditioning is removed. The current maximum-height proof alone does not
establish such a point-process limit: multiple columns in one unimodal unit
can be dependent and must be controlled explicitly.

\subsection{The first correction to the mixed extreme law}
The present error in \eqref{vmb:cw:eq:mixture} is $O(\log n/n)$. Derive the full
correction, including the perturbed residues, the next kernel-tail term,
and the simplex boundary terms in the coefficient sum. The exact
$\Psi$-expansion and the cutoff polynomials show that the correction should
involve the lattice phase and powers of $\log n$. This would sharpen
critical height quantiles beyond the current rounding band.

\subsection{Growing numbers of critical levels}
All multicritical theorems here keep $m$ fixed. Determine how fast $m=m(n)$
may grow while the Dirichlet approximation remains uniform. Since
$\rho_{m+1}\to1$, analytic neighborhoods shrink and the number of possible
competing poles increases. Neither the fixed-$m$ simplex proof nor the
fixed-minimum-part asymptotic theorem justifies this interchange of limits.

\subsection{A local limit theorem for coexistence fractions}
The bounded-support weak limits proved here determine macroscopic moments.
A quantitative joint local limit for $(V_{1,n},\ldots,V_{m,n})$ would describe
individual lattice probabilities, including faces of the simplex. Interior
Gaussian fluctuations around a prescribed Dirichlet fraction should interact
with boundary layers where one renewal family has only finitely many units.
A uniform theorem across those layers remains to be proved.

\subsection{Nonreal poles and optimal exponential errors}
The proofs require a radius slightly beyond the dominant singularity, not a
certification of the next singularity. Locate and certify the first several
nonreal poles of the weighted and height-capped models. This would make the
exponential remainder effective, identify oscillatory correction sectors,
and improve the discrete counting inverses near difficult thresholds.

\subsection{Natural boundaries and moving minimum parts}
The unit-circle natural-boundary question and the moving-lower-bound crossover
problem in Part~I remain separate questions. The present $q$-product
identity may help study root-of-unity behavior, but it does not prove that
singularities are noncancelling near every boundary point. Likewise, increasing
height cutoff and increasing minimum part are different limits; the
maximum-height theorem does not solve the latter crossover problem.

\begin{writenote}[Added 5 October 2026, batch~98]
These are Part~I's questions ``Is the unit circle a natural boundary?'' and
``Moving lower bounds and a crossover problem'' (Section~\ref{vmb:sec:future}),
restated; Part~II answers neither.
\end{writenote}

\subsection{Formalization and effective certificates}
\label{vmb:cw:sec:q-formal}
A first formalization target is the unique marked decomposition, including
plateau semantics and height caps. Next come exact truncated-series identities
and the elementary simplex coefficient lemma. Certified interval bounds for
$u_c$, $\alpha$, and the critical intensities would be useful independently
of a full complex-analytic formalization. The numerical agreement in this
package is not a replacement for those certificates.

\subsection{Sketched claims of the manuscript, and what remains of Part~I's question}
\label{vmb:cw:sec:q-write}
\begin{writenote}[Added 5 October 2026, batch~98]
Each item names the manuscript's claim, its support there, and its status
now. The proofs supplied in writing have not been independently reviewed;
until they are, the items they settle should be read as settled only to that
standard.
\begin{enumerate}[leftmargin=*,itemsep=2pt]
\item \emph{Conditional intensity of the critical maximum} (manuscript
 Section~8.1; here the paragraph before
 Proposition~\ref{vmb:cw:prop:conditional}). Claim: conditionally on the limiting
 area fractions $X_i$, the extreme-height law has intensity
 $\sum\alpha_i^*X_i$, so the critical maximum is generally not
 asymptotically independent of the fractions. Support: one sentence (``Joint
 marking of the area fractions shows more''). Missing: the simplex lemma with
 simultaneous area marks and the cutoff, and the identification of the joint
 limit. Status: proved in writing, Proposition~\ref{vmb:cw:prop:conditional},
 which also makes ``generally'' exact (independence holds if and only if all
 $\alpha_i^*$ coincide).
\item \emph{Tilted multicritical window} (manuscript Section~5.2). Claim: the
 weights $u_j^*e^{s_j/n}$ multiply the leading coefficient by
 $\mathcal G_m(\bm s)$ and tilt the Dirichlet law. Support: ``The proof is the
 same nearby-pole argument.'' Status: proved in writing,
 Proposition~\ref{vmb:cw:prop:tilted}.
\item \emph{Complete fixed-parameter classification} (manuscript
 Section~5.2). Claim: with $\nu$ minimal roots, the pole has order $\nu$ and
 the coefficients are $r_*^{-n}$ times a polynomial of degree $\nu-1$.
 Support: a paragraph sketch (noncancellation through the suffix after the
 last active level). Status: proved in writing,
 Proposition~\ref{vmb:cw:prop:chambers}.
\item \emph{Multicritical rounding and critical quantiles} (manuscript
 Sections~9.3 and~9.4). Claims: rounding bounds ``of the form''
 \eqref{vmb:cw:eq:rounding} at every multicritical point, and the critical
 analogue of \eqref{vmb:cw:eq:quantilerounding} through the unique crossing of
 the mixture transform. Support: analogy. Status: proved in writing,
 Proposition~\ref{vmb:cw:prop:roundingmulti}.
\item \emph{Monotonicity of the interpolated cutoff shift} (manuscript
 Section~9.1). Claim: $\delta(h)$ is strictly decreasing for large real $h$.
 Support: one line. Status: proved in writing,
 Lemma~\ref{vmb:cw:lem:deltamono}.
\item \emph{The residue ratio} \eqref{vmb:cw:eq:residueconv} in the proof of
 Theorem~\ref{vmb:cw:thm:gumbel}: argued in outline; recorded as complete in
 outline, not as an open question.
\item \emph{Part~I's question ``Joint limits and extreme heights''}
 (Section~\ref{vmb:sec:future}): its extreme-height half is answered
 (Theorems~\ref{vmb:cw:thm:gumbel}, \ref{vmb:cw:thm:independent}). Its other
 half stays open and is not addressed by the manuscript: a joint normal law
 for the additive statistics that the renewal representation can mark (number
 of columns, number of ones, peak heights), which needs a proof that the
 relevant covariance matrix is nondegenerate.
\end{enumerate}
No claim of the manuscript was found to be wrong.
\end{writenote}

\section{Claim ledger and conclusion of Part II}
\label{vmb:cw:sec:ledger}
\begin{longtable}{@{}>{\raggedright\arraybackslash}p{0.35\textwidth}>{\raggedright\arraybackslash}p{0.59\textwidth}@{}}
\toprule
Statement & Status and boundary\\
\midrule
\endfirsthead
\toprule
Statement & Status and boundary (continued)\\
\midrule
\endhead
A001523 formulas and classical asymptotics & Prior literature and OEIS;
not claimed as new. Only the block connection is used.\\[3pt]
Minimum-part decomposition and unweighted simple pole & Prior results,
reproved here to establish the hypotheses.\\[3pt]
Bottom-valley phase transition & Proved for fixed positive weight; exact
critical value, three regimes, and critical window are explicit.\\[3pt]
All window orders & Finite-order, uniform additive expansions, with an
algorithm from local Taylor and Laurent data.\\[3pt]
Multicritical Dirichlet law & Proved for any fixed finite number of marked
levels; no growing-level uniformity is asserted.\\[3pt]
Ordinary maximum-height law & Proves the lattice-Gumbel question posed in
the previous report, with explicit intensity and $O(\log n/n)$ window error.\\[3pt]
Critical maximum-height law & Proved Dirichlet-mixture formula, including
simultaneous $1/n$ weight tilts and a uniform height-window error.\\[3pt]
Cutoff-pole expansion and inverse & All finite exponential orders with
polynomial coefficients and explicit remainders; no unrestricted convergence
or Stokes claim.\\[3pt]
Discrete inverses and quantiles & Specified smooth interpolants plus two
ceilings; no unconditional rounding shortcut.\\[3pt]
Exact tests and numerical data & Independent finite integer checks pass;
decimal values are reproducible diagnostics, not interval-certified constants.\\[3pt]
Novelty and formal status & An extension of documented repository work;
no exhaustive priority claim, no Lean verification, and no independent peer
review are asserted.\\
\bottomrule
\end{longtable}

\begin{writenote}[Added 5 October 2026, batch~98]
``The previous report'' is Part~I. The row ``Cutoff-pole expansion and
inverse'' was delivered as ``Cutoff-pole transseries and inverse''
(see the note after Theorem~\ref{vmb:cw:thm:cutoffseries}); the table is set
ragged-right. The statements
completed in writing are listed in Subsection~\ref{vmb:cw:sec:q-write}.
\end{writenote}

The central conclusion is that the same OEIS-linked renewal decomposition
supports two different kinds of randomness. In the ordinary unweighted
regime, one renewal family occupies almost all the area, its valley count is Gaussian,
and its tallest column obeys a lattice-Gumbel law independent of that count.
At coexistence, several families share macroscopic area with a Dirichlet
law; their joint extreme distribution retains that random allocation as a
mixture intensity. The pole multiplicity controls the counting polynomial,
the critical-window scale, the macroscopic covariance, and the appropriate
Lambert-$W$ inverse. The exact height-tail identity makes the exponentially
small cutoff corrections computable rather than merely qualitative.

\clearpage
\phantomsection\addcontentsline{toc}{part}{Appendices of Part I}
\appendix
\section{A coefficient recurrence suitable for exact checking}
Let all series be truncated modulo~\(q^{N+1}\). For descending
\(j=N,N-1,\ldots,1\), compute
\begin{align}
Q_j&=1-q^j-q^jU_{j+1},\notag\\
(1-q^j)^2U_j&=q^j(1-q^j)+U_{j+1},\label{vmb:eq:algorithm}\\
(1-q^j)Q_jD_j&=q^jQ_j+D_{j+1}.\notag
\end{align}
Every denominator has constant term one. If
\(v(q)w(q)=b(q)\), \(v(0)=1\), its coefficients are obtained by
\[
w_n=b_n-\sum_{k=1}^n v_k w_{n-k}.
\]
This requires no division in the coefficient ring. Initialization with
zero at level \(N+1\) is exact at the requested order, so integrality
of every computed coefficient is automatic. The enumeration proof in
Section~\ref{vmb:sec:decomposition}, not integrality alone, identifies the
objects being counted.

\section{Claim ledger and limitations}
\begin{longtable}{@{}p{0.33\textwidth}p{0.60\textwidth}@{}}
\toprule
Claim & Status and justification\\
\midrule
\endfirsthead
\toprule
Claim & Status and justification (continued)\\
\midrule
\endhead
Published generating function & Prior result of~\cite{FRV}; independently
rederived by the unique minimum-part decomposition.\\[3pt]
Pure exponential equivalent & Proved in Theorem~\ref{vmb:thm:leading}, with
exact constants and an exponentially smaller error.\\[3pt]
Corrected numerical constants & Certified by standard-library integer
interval arithmetic with explicit infinite-tail bounds.\\[3pt]
Infinite positive poles; non-D-finiteness & Proved without a
noncancellation assumption in Theorem~\ref{vmb:thm:ladder} and
Corollary~\ref{vmb:cor:nonD}.\\[3pt]
Arbitrary exponential cutoff & Proved in Theorem~\ref{vmb:thm:sectors}; not
a claim of convergence of an infinite pole sum.\\[3pt]
Limit laws & Proved by uniform analytic perturbation and direct
characteristic-function convergence.\\[3pt]
All logarithmic pole orders & Proved in Theorem~\ref{vmb:thm:W}; the
\(1/j\)-scale correction remains to be sharpened.\\[3pt]
Inverse expansion & Proved for a specified finite-sector interpolant;
integer crossings are bounded with two ceilings.\\[3pt]
Novelty and formal status & No matching prior proof found in the targeted
search. No exhaustive priority claim, no Lean proof, and no independent
peer review are asserted.\\
\bottomrule
\end{longtable}

\clearpage
\phantomsection\addcontentsline{toc}{part}{Appendices of Part II}
\section{Computing the Laurent constants without unstable subtraction}
\label{vmb:cw:app:laurent}
At $\rho_2$, write
\[
 D_2(q)=\frac{c_2}{1-q/\rho_2}+c_{2,0}+O(1-q/\rho_2).
\]
Let
\[
 B_2(q)=E_2(q)D_3(q),\quad
 M_2(\rho_2)=\rho_2K_2'(\rho_2),\quad J_2=\rho_2^2K_2''(\rho_2)/2.
\]
Since $K_2(\rho_2)=1$, direct expansion of \eqref{vmb:cw:eq:Drec} gives
\begin{align}
 c_2&=B_2(\rho_2)/M_2(\rho_2),\\
 c_{2,0}&=H_2(\rho_2)+\frac{B_2(\rho_2)J_2}{M_2(\rho_2)^2}-\frac{\rho_2B_2'(\rho_2)}{M_2(\rho_2)},\\
 b&=\frac{c_2}{(1-\rho_2)^2},\\
 b_0&=\frac{c_{2,0}}{(1-\rho_2)^2}-\frac{2\rho_2c_2}{(1-\rho_2)^3}.
\end{align}
These formulas avoid subtracting two large nearby singular values of $B$.
They are the formulas used in the numerical program.

\begin{writenote}[Added 5 October 2026, batch~98]
Here $c_2$ is the dominant amplitude of $D_2$ (as in
Theorem~\ref{vmb:thm:leading}), not the amplitude $C_2$ of $D_1$ at $\rho_2$;
the manuscript's $A_2$, $d_{20}$ and $M_2=RK_2'(R)$ are printed as $B_2$
($B_2=E_2D_3$ is the level-2 analogue of $B=E_1D_2$), $c_{2,0}$ and
$M_2(\rho_2)$.
\end{writenote}

\section{A fully finite prescription for all height-cutoff coefficients}
\label{vmb:cw:app:cutoff}
For a requested order $N$, compute $\gamma_0(q),\ldots,\gamma_{N-1}(q)$ by
\eqref{vmb:cw:eq:qbinomcoeff} and then the first $N$ coefficients of $\Psi$ by
\eqref{vmb:cw:eq:psicoeff}. Put
\[
 A_r(q)=H_1(q)P_2(q)\,[z^r]\Psi(q,z)\,q^r.
\]
These satisfy $K(q)-K^{\le h}(q)=\sum_{r\ge1}A_r(q)q^{rh}$.
Now introduce a formal indeterminate $\epsilon$ and polynomials
$\varpi_0(h),\ldots,\varpi_{N-1}(h)$. In
\[
 K\!\left(\rho\exp\left(\sum_{r=1}^N\varpi_{r-1}(h)\epsilon^r\right)\right)-1
 =\sum_{r=1}^N A_r\!\left(\rho\exp\left(\sum_{\ell=1}^N
              \varpi_{\ell-1}(h)\epsilon^\ell\right)\right)
 \epsilon^r\exp\left(rh\sum_{\ell=1}^N\varpi_{\ell-1}(h)\epsilon^\ell\right),
\]
compare coefficients through $\epsilon^N$. At step $r$ divide the known
residual by $M_\rho$ to obtain $\varpi_{r-1}$. Only derivatives of analytic
functions at one ordinary point $\rho<1$ occur. Theorem
\ref{vmb:cw:thm:cutoffseries} proves that this finite symbolic prescription has
the asserted asymptotic meaning.

\begingroup
\raggedright
\begin{thebibliography}{9}
\bibitem{FRV}
R. Fl\'orez, J. L. Ram\'irez, and D. Villamizar,
\emph{Restricted bargraphs and unimodal compositions},
Journal of Combinatorial Theory, Series A \textbf{208} (2024), 105934.
\href{https://doi.org/10.1016/j.jcta.2024.105934}{doi:10.1016/j.jcta.2024.105934}.
Especially Section~2 and Conjecture~2.3. Publisher HTML consulted
October~1, 2026.

\bibitem{OEIS}
OEIS Foundation Inc., \emph{A001523: Number of stacks, or planar
partitions of \(n\); also number of weakly unimodal compositions of
\(n\)}. \url{https://oeis.org/A001523}.
Entry consulted October~1, 2026. The entry credits Auluck's known leading
asymptotic and links the bargraph paper above.

\bibitem{FS}
P. Flajolet and R. Sedgewick, \emph{Analytic Combinatorics},
Cambridge University Press, 2009. Chapters~IV--V for rational and
meromorphic asymptotics, and Chapter~IX for parameter distributions.
Author's book page: \url{https://sedgewick.io/books/analytic-combinatorics/}.

\bibitem{ProveIt}
V. Reshetnikov, \emph{ProveIt}, project repository, particularly
\texttt{Analysis/Transseries/README.md}.
\url{https://github.com/VladimirReshetnikov/ProveIt}.
Read-only inspection, October~1, 2026. No repository changes were made.

\bibitem{vmb-tai}
\emph{Transseries and inversion}, a volume of the ProveIt repository:
\nolinkurl{Analysis/Transseries/docs/series-and-transseries/Transseries_And_Inversion/transseries_and_inversion.tex}
(theorems \texttt{p0:thm:lambert-core}, \texttt{p0:thm:core-reversion},
\texttt{p0:thm:staircase} and \texttt{p0:thm:LB}). Added in writing,
batch~75.
\end{thebibliography}
\endgroup
\end{document}
